\documentclass[12pt,a4paper,UTF8]{ctexart}
\usepackage[a4paper,top=22mm,bottom=22mm,left=22mm,right=22mm,headheight=15pt]{geometry}
\newcommand{\Needspace}[1]{\par\begingroup\dimen0=#1\relax\dimen2=\pagegoal\advance\dimen2 by-\pagetotal\ifdim\dimen0>\dimen2\newpage\fi\endgroup}
\newcommand{\needspace}[1]{\par\begingroup\skip0=\lastskip\vskip-\skip0\dimen0=#1\relax\advance\dimen0 by\skip0\vskip\dimen0\penalty-100\vskip-\dimen0\vskip\skip0\penalty9999\endgroup}
\usepackage{amsmath,amssymb,mathtools,bm}
\usepackage{graphicx,xcolor,booktabs,longtable,array,enumitem,fancyhdr,caption,float,microtype,setspace}
\usepackage[hidelinks,unicode]{hyperref}
\usepackage[most]{tcolorbox}
\definecolor{ContestBlue}{HTML}{17365D}
\hypersetup{pdftitle={第064套},pdfauthor={}}
\setstretch{1.24}
\setlength{\parindent}{2em}
\setlength{\parskip}{0.35em}
\setlength{\emergencystretch}{2em}
\setlist{itemsep=0.35em,topsep=0.4em}
\captionsetup{font=small,labelfont=bf,justification=centering}
\setcounter{secnumdepth}{0}
\pagestyle{fancy}
\fancyhf{}
\fancyhead[L]{\small 第064套}
\fancyfoot[C]{\thepage}
\renewcommand{\headrulewidth}{0.4pt}
\ctexset{
  section={format=\Large\bfseries\color{ContestBlue},beforeskip=1.2em,afterskip=0.7em},
  subsection={format=\large\bfseries\color{ContestBlue},beforeskip=1.0em,afterskip=0.5em},
  subsubsection={format=\normalsize\bfseries,beforeskip=0.8em,afterskip=0.35em}
}
\usepackage{calc}
\newcounter{none}
\let\originalsubsection\subsection
\renewcommand{\subsection}{\needspace{6\baselineskip}\originalsubsection}
\makeatletter
\newcommand{\tightlist}{\setlength{\itemsep}{0pt}\setlength{\parskip}{0pt}}
\newcommand{\maxwidth}{\ifdim\Gin@nat@width>\linewidth\linewidth\else\Gin@nat@width\fi}
\newcommand{\maxheight}{\ifdim\Gin@nat@height>0.72\textheight 0.72\textheight\else\Gin@nat@height\fi}
\newcommand{\pandocbounded}[1]{\begingroup\setkeys{Gin}{width=\maxwidth,height=\maxheight,keepaspectratio}#1\endgroup}
\makeatother
\begin{document}
\begin{titlepage}
  \centering
  \vspace*{58mm}
  {\Huge\bfseries\color{ContestBlue} 第064套\par}
  \vfill
  \vspace*{28mm}
\end{titlepage}
\section{第 1 题}\label{ux7b2c-1-ux9898}

如图，\(A\)、\(B\)、\(C\) 三个物体的质量均为 \(m\)。\(A\) 与
\(C\)、\(B\) 与 \(C\) 之间的摩擦因数分别为 \(\mu_1\)、\(\mu_2\)，\(C\)
与地面之间的摩擦因数为 \(\mu_3\)。轻绳绕过安装在 \(A\)、\(B\)
上的光滑滑轮，绳的自由端受水平拉力 \(F\)，且 \(A\)、\(B\) 相距足够远。

（1）从 \(F=0\) 开始缓慢增大拉力，直至有物体开始运动。以
\(x=\mu_1/\mu_3\) 为横轴、\(y=\mu_2/\mu_3\)
为纵轴，划分最先发生各种运动的参数区域，并给出边界方程。

（2）在首次运动发生后继续增大
\(F\)，讨论各摩擦因数组合下可能出现的运动状态。只需说明哪些物体运动以及
\(C\) 的运动方向。

\begin{figure}[H]
\centering
\includegraphics[width=0.76\linewidth,height=\textheight,keepaspectratio,alt={图1-1　三物体、滑轮与轻绳的连接关系}]{../figures/question-01.pdf}
\caption*{图1-1　三物体、滑轮与轻绳的连接关系}
\end{figure}

\subsection{一、题型定位}\label{ux4e00ux9898ux578bux5b9aux4f4d}

先沿绳路数清张力：\(A\) 受 \(3F\) 向右，\(B\) 受 \(2F\)
向左。首次启动由三处静摩擦临界的先后决定；启动以后，还要区分相对滑动、对地运动和加速度方向。

\subsection{二、知识点}\label{ux4e8cux77e5ux8bc6ux70b9}

\begin{enumerate}
\def\labelenumi{\arabic{enumi}.}
\tightlist
\item
  理想轻绳和光滑滑轮中，同一根绳上的张力处处为 \(F\)。
\item
  动滑轮受到相邻两段绳张力的矢量合力；结合绳端受力后，\(A\)、\(B\) 对
  \(C\) 的相对滑动趋势分别对应 \(3F\)、\(2F\) 的切向需求。
\item
  静摩擦在临界前是被动力，必须写成不等式
  \(|f_s|\leq\mu N\)；只有临界或滑动后才能取极限值。
\item
  \(C\) 对地面的正压力为三物体总重 \(3mg\)。
\item
  参数分区边界由 \(F_A=F_B\)、\(F_A=F_C\)、\(F_B=F_C\) 给出。
\item
  两个上接触面均滑动后，\(C\) 的受力方向由 \(\mu_1-\mu_2\) 的符号决定。
\end{enumerate}

\subsection{三、物理图像}\label{ux4e09ux7269ux7406ux56feux50cf}

图 1-1 中绳端拉力为 \(F\)，理想绳张力也为 \(F\)。沿绳路逐段观察，安装在
\(A\) 上的滑轮及绳端共同使 \(A\) 相对 \(C\) 的水平驱动力大小为
\(3F\)；安装在 \(B\) 上的滑轮受到两段张力，合力大小为
\(2F\)。因此若系统仍静止，上表面两处所需静摩擦力的大小分别是

\[
f_{1,\mathrm{req}}=3F,\qquad f_{2,\mathrm{req}}=2F.
\]

把 \(A+B+C\)
看成整体，绳与滑轮之间的力均为内力，系统所受水平外力只有绳自由端的 \(F\)
与地面对 \(C\) 的摩擦力，所以

\[
f_{3,\mathrm{req}}=F.
\]

这三个``所需摩擦力''与三个极限值比较，便得到第一问全部边界。

\subsection{四、解题入口}\label{ux56dbux89e3ux9898ux5165ux53e3}

先不要令摩擦力等于 \(\mu N\)。在静止阶段依次写出

\[
3F\leq\mu_1mg,\qquad
2F\leq\mu_2mg,\qquad
F\leq 3\mu_3mg.
\]

于是三处接触面单独达到临界时的外力为

\[
F_A=\frac{\mu_1mg}{3},\qquad
F_B=\frac{\mu_2mg}{2},\qquad
F_C=3\mu_3mg.
\]

最小者决定谁先动。再把两两相等条件除以
\(\mu_3\)，就得到参数平面中的三条边界。

\subsection{五、解析}\label{ux4e94ux89e3ux6790}

\begin{figure}[H]
\centering
\includegraphics[width=0.9\linewidth,height=\textheight,keepaspectratio,alt={图1-2　三处接触面的所需摩擦力与临界比较}]{../figures/solution-01.pdf}
\caption*{图1-2　三处接触面的所需摩擦力与临界比较}
\end{figure}

图1-2左侧只保留首次启动的有效边界段，右侧另画两处上层接触均滑动时的受力带。底块已运动时需要速度历史，不能用受力带直接判定立即停下。

图 1-2 将 \(A\)、\(B\) 和整体分别隔离。图中的 \(3F\)、\(2F\)、\(F\)
是由同一绳张力沿实际绳路合成后得到的三个切向需求，后续分区与第二阶段受力均使用这些量。

\needspace{3\baselineskip}
\textbf{（一）静止阶段的三个临界外力}

\(A\) 与 \(C\) 之间的正压力为

\[
N_1=mg.
\]

保持相对静止所需的摩擦力大小为 \(3F\)，所以静止条件为

\[
3F\leq \mu_1mg.
\]

等号成立时 \(A\) 相对 \(C\) 达到滑动临界：

\[
F_A=\frac{\mu_1mg}{3}. \tag{1}
\]

同理，\(B\) 与 \(C\) 之间有 \(N_2=mg\)，所需摩擦力为 \(2F\)：

\[
2F\leq\mu_2mg,
\]

故

\[
F_B=\frac{\mu_2mg}{2}. \tag{2}
\]

地面承担三个物体的总重，因此

\[
N_3=3mg.
\]

把三个物体和滑轮、轻绳视为整体，静止时地面摩擦力平衡水平外力 \(F\)：

\[
f_3=F.
\]

地面临界条件为

\[
F=\mu_3N_3=3\mu_3mg,
\]

即

\[
F_C=3\mu_3mg. \tag{3}
\]

式（1）---（3）的量纲均为力，这是第一层检查；当任一摩擦因数趋于零时，对应临界外力也趋于零，极限行为正确。

\needspace{3\baselineskip}
\textbf{（二）两两临界相等所给出的边界}

定义

\[
x=\frac{\mu_1}{\mu_3},\qquad
y=\frac{\mu_2}{\mu_3}.
\]

令 \(F_A=F_B\)：

\[
\frac{\mu_1mg}{3}=\frac{\mu_2mg}{2}.
\]

约去公共因子 \(mg\)，交叉相乘：

\[
2\mu_1=3\mu_2.
\]

两边除以 \(3\mu_3\)：

\[
\boxed{y=\frac{2}{3}x}. \tag{4}
\]

令 \(F_A=F_C\)：

\[
\frac{\mu_1mg}{3}=3\mu_3mg.
\]

约去 \(mg\)，再乘以 \(3\)：

\[
\mu_1=9\mu_3,
\]

所以

\[
\boxed{x=9}. \tag{5}
\]

令 \(F_B=F_C\)：

\[
\frac{\mu_2mg}{2}=3\mu_3mg,
\]

得到

\[
\mu_2=6\mu_3,
\]

所以

\[
\boxed{y=6}. \tag{6}
\]

三条边界交于

\[
(x,y)=(9,6),
\]

且式（4）在 \(x=9\) 时确有 \(y=6\)，说明三处临界可同时出现。

\needspace{3\baselineskip}
\textbf{（三）最先运动区域}

\begin{enumerate}
\def\labelenumi{\arabic{enumi}.}
\tightlist
\item
  \(A\) 先相对 \(C\) 滑动要求 \(F_A<F_B\) 且 \(F_A<F_C\)。
\end{enumerate}

由 \(F_A<F_B\)：

\[
\frac{\mu_1}{3}<\frac{\mu_2}{2}
\Longleftrightarrow
y>\frac23x.
\]

由 \(F_A<F_C\)：

\[
\frac{\mu_1}{3}<3\mu_3
\Longleftrightarrow x<9.
\]

因此

\[
\boxed{A\text{ 先滑： }x<9,\quad y>\frac23x}. \tag{7}
\]

\begin{enumerate}
\def\labelenumi{\arabic{enumi}.}
\setcounter{enumi}{1}
\tightlist
\item
  \(B\) 先相对 \(C\) 滑动要求 \(F_B<F_A\) 且 \(F_B<F_C\)：
\end{enumerate}

\[
\boxed{B\text{ 先滑： }y<6,\quad y<\frac23x}. \tag{8}
\]

\begin{enumerate}
\def\labelenumi{\arabic{enumi}.}
\setcounter{enumi}{2}
\tightlist
\item
  \(C\) 先相对地面滑动要求 \(F_C<F_A\) 且 \(F_C<F_B\)：
\end{enumerate}

\[
\boxed{C\text{ 先滑： }x>9,\quad y>6}. \tag{9}
\]

首次启动的有效边界只有三段：\(y=2x/3\) 且 \(0\le x\le9\)；\(x=9\) 且
\(y\ge6\)；\(y=6\) 且
\(x\ge9\)。其余延长线上虽有两个临界力相等，却已有第三处先滑，不能画作首次启动的区域边界。交点
\((9,6)\) 为三处同时临界。这里要求
\(\mu_3>0\)；无地面摩擦的情形须直接比较原始力，不能除以 \(\mu_3\)。

\needspace{3\baselineskip}
\textbf{（四）启动后的中间状态}

以下按题目的单一摩擦因数模型，取最大静摩擦与滑动摩擦均为
\(\mu_jN_j\)。令 \(f=F/(\mu_3mg)\)，各加速度以 \(\mu_3g\) 为单位。用
\(f_1,f_2\) 表示 \(C\) 从上层受到的向右、向左摩擦力，力也按 \(\mu_3mg\)
归一化。牛顿方程为

\[
a_A=3f-f_1,\qquad a_B=-2f+f_2,\qquad a_C=f_1-f_2-f_3.
\]

静摩擦满足
\(|f_1|\le x,|f_2|\le y,|f_3|\le3\)；摩擦方向须由当前相对速度确定。刚开始通常有以下三种状态。

\begin{itemize}
\tightlist
\item
  \(A\) 先向右滑、\(B,C\) 暂时静止：\(f_1=x,f_2=2f\)，所以 \(C\)
  仍静止的条件是 \(|x-2f|\le3\)，且 \(2f\le y\)。
\item
  \(B\) 先向左滑、\(A,C\) 暂时静止：\(f_1=3f,f_2=y\)，条件是
  \(|3f-y|\le3\)，且 \(3f\le x\)。
\item
  地面先滑、三块暂时无相对滑动：令 \(d=\operatorname{sgn}v_C\)，则
  \(a=(f-3d)/3\)，接触摩擦为
  \(f_1=3f-a,f_2=2f+a\)，须仍满足两个上层静摩擦上限。首次地面启动时
  \(d=1\)。
\end{itemize}

当其中某一不等式达到边界，重新检查接触状态。例如 \(A\) 已滑而 \(B,C\)
对地一起运动时，

\[
a_{BC}=\frac{x-2f-3d}{2},\qquad f_2=2f+a_{BC},\qquad |f_2|\le y.
\]

若 \(B\) 已滑而 \(A,C\) 一起运动，则

\[
a_{AC}=\frac{3f-y-3d}{2},\qquad f_1=3f-a_{AC},\qquad |f_1|\le x.
\]

这些式子把``只有上层滑''\,``整组滑''\,``上层滑且底块也动''的中间阶段补齐。不能用大外力下的最终分类代替启动后的全部过程。

\needspace{3\baselineskip}
\textbf{（五）两处上接触均滑后的方向}

在 \(A\) 相对 \(C\) 向右、\(B\) 相对 \(C\)
向左滑动时，\(f_1=x,f_2=y\)。若 \(C\) 尚静止，地面可保持静止当且仅当

\[
\boxed{|x-y|\le3}.
\]

若 \(C\) 已在运动，则应使用滑动摩擦：

\[
\boxed{a_C=x-y-3\operatorname{sgn}v_C}.
\]

因此 \(x-y>3\) 对应最终向右加速，\(y-x>3\) 对应最终向左加速。\(|x-y|<3\)
时，已在运动的 \(C\)
会先减速，再停住；不是一进入此区就立即静止。等号时同向运动的 \(C\)
可保持非零恒速，而静止的 \(C\) 也能处于临界。

第二问没有给出 \(F(t)\) 和各阶段持续时间，因而仅凭 \(x,y\)
不能确定每一时刻的对地速度。上式和接触不等式给出可判定的运动状态；若要画完整时间轨迹，必须再给拉力的时间变化。图1-2分别画首次启动的有效边界和两处均滑时的受力分类，后者不作为完整运动历史图。

\subsection{六、关键洞察}\label{ux516dux5173ux952eux6d1eux5bdf}

解析图左侧只保留首次启动的有效边界段，右侧另画两处上层接触均滑动时的受力带。底块已运动时需要速度历史，不能用受力带直接判定立即停下。

第一问不是求一个临界，而是比较三个临界。参数平面中的每条直线都应能追溯到``哪两个临界外力相等''。第二问则换了比较对象：上层两处已经滑动，摩擦大小不再随
\(F\) 增长，\(C\) 是否运动只取决于二者之差能否被地面静摩擦抵消，因此出现
\(|x-y|=3\)。

参数图中，\(x=9\)、\(y=6\) 决定地面与两上层临界的先后，斜线 \(y=2x/3\)
决定两上层彼此的先后；而第二问的两条平行线 \(y=x\pm3\)
决定滑动后地面能否继续保持静止。前一组边界来自``比值''，后一组边界来自``差值''，二者不能互相替代。两组线分别描述首次启动和两处上层均滑后的受力条件；中间状态还须应用解析中的动力学方程。

\subsection{七、易错点}\label{ux4e03ux6613ux9519ux70b9}

\begin{enumerate}
\def\labelenumi{\arabic{enumi}.}
\tightlist
\item
  把滑轮当作固定在地面的定滑轮，漏掉滑轮支架传给 \(A\)、\(B\) 的合力。
\item
  在系统尚未临界时直接写 \(f_1=\mu_1mg\)、\(f_2=\mu_2mg\)。
\item
  忘记地面正压力是 \(3mg\)，把 \(F_C\) 错写为 \(\mu_3mg\)。
\item
  只比较 \(F_A\) 与 \(F_B\)，遗漏地面先滑区域 \(x>9,y>6\)。
\item
  把边界 \(y=2x/3\) 写成 \(y=3x/2\)。可用交点 \((9,6)\) 立即核查斜率。
\item
  第二问把两个上层摩擦力相加；它们对 \(C\) 的方向相反，应取差。
\item
  只说``\(C\) 会动''，没有用 \(x-y\) 的符号判断其向左还是向右。
\item
  混淆``相对 \(C\) 滑动''和``相对地面运动''。
\end{enumerate}

批改时建议按四项核验：绳路倍数是否正确、临界前是否使用静摩擦不等式、三种首次启动临界是否全部比较、第二问是否用摩擦力之差判断
\(C\)。只要第一项错，后面的参数图即使画得整齐也不能给主要步骤分；若前三项正确而只漏边界等号，可视为临界状态表述不完整。

\subsection{八、课堂讲解建议}\label{ux516bux8bfeux5802ux8bb2ux89e3ux5efaux8bae}

先遮住摩擦系数，只让学生沿实际绳路在每一段标出张力 \(F\)，分别圈出``物体
\(A\) 加其滑轮''\,``物体 \(B\) 加其滑轮''，算出水平合力。确认
\(3F\)、\(2F\)、\(F\) 后，再让学生把三个临界外力写在同一行并排序。

画参数图时按三步进行：先画斜线 \(y=2x/3\)，再画 \(x=9\) 与
\(y=6\)，最后用测试点判断区域。第二问可用三个数值点核验：\((x,y)=(7,3)\)
时 \(C\) 向右；\((3,7)\) 时 \(C\) 向左；\((4,5)\) 时静摩擦能够保持 \(C\)
静止，已动的 \(C\) 先减速。

课堂板书建议保留``静止阶段''和``滑动阶段''两栏。左栏只写
\(f_{\rm req}\leq f_s^{\max}\)，右栏才写
\(f_k=\mu N\)，借此让学生看见同一个摩擦系数在两个阶段扮演的角色不同。比较临界力时，可先让学生口头回答``若
\(x\) 很小而 \(y\)
很大，谁最容易滑''，再把这一极端判断与参数图对应，防止只会机械画线。

讲完后可做一次反向诊断：给出区域条件 \(x>9,y>6\)，要求学生不用图直接比较
\(F_A,F_B,F_C\)；再给出 \(|x-y|<3\)，要求说明为何这时 \(C\) 不动但
\(A,B\)
可以相对它滑动。若学生把两个问题混在一起，说明他还没有分清``首次启动''和``继续增大外力后的状态''。

\subsection{九、变式训练}\label{ux4e5dux53d8ux5f0fux8badux7ec3}

\begin{enumerate}
\def\labelenumi{\arabic{enumi}.}
\tightlist
\item
  若 \(A\)、\(B\)、\(C\) 的质量分别为
  \(m_A,m_B,m_C\)，重新求三条第一临界边界。
\item
  若地面只能提供大小固定的滑动摩擦，讨论 \(C\) 启动后的加速度方向。
\item
  把绳端拉力改为恒定加速度拉动的质量块，重新建立绳长约束与临界条件。
\item
  保持 \(x\) 不变逐渐增大 \(y\)，在参数图上描述最先运动物体和最终 \(C\)
  运动方向的变化。
\end{enumerate}

\Needspace{28\baselineskip}

\textbf{变式解析要点：} 对第 1 项，绳路给出的 \(3F\)、\(2F\)
不变，但极限摩擦改为 \(\mu_1m_Ag\)、\(\mu_2m_Bg\)，地面支持力改为
\((m_A+m_B+m_C)g\)，因此先写

\[
F_A=\frac{\mu_1m_Ag}{3},\quad
F_B=\frac{\mu_2m_Bg}{2},\quad
F_C=\mu_3(m_A+m_B+m_C)g,
\]

再两两相等即可。第 2 项启动后若三层仍接触，\(C\) 的加速度方向由

\[
\mu_1m_Ag-\mu_2m_Bg-f_3
\]

的符号决定；滑动摩擦 \(f_3\) 必与 \(C\) 的速度相反，不能预先固定方向。第
4 项沿竖线 \(x=\text{常量}\) 上移，依次记录穿过 \(y=2x/3\)、\(y=6\) 与
\(y=x\pm3\) 的顺序，便能把``首次启动''和``后续方向''放在同一张参数图上。

\newpage

\section{第 2 题}\label{ux7b2c-2-ux9898}

在光滑的一个抛物面 \(z=\dfrac a2(x^2+y^2)\) 上放一个质量为 \(m\)
的质点，重力加速度 \(g\) 沿着 \(z\)
轴负方向。任意给定初态后质点开始振动，振动的幅度远小于 \(1/a\)，其中
\(a>0\)。

（1）计算两个方向独立的简谐振动的角频率 \(\omega_0\)。

（2）换到一个以恒定角速度 \(\omega_1\) 绕 \(z\)
轴转动的参照系中，新的水平坐标轴记为 \(X',Y'\)，写下 \(X'\)、\(Y'\)
应当满足的微分方程。

（3）解出（2）中的微分方程，并解释这个结果如何和（1）中的结果对应。

（4）我们把抛物面改为马鞍面
\(z=\dfrac a2x^2-\dfrac a4y^2\)，这样显然振动是不稳定的，但是如果让这个面绕着
\(z\) 轴以恒定的角速度 \(\omega_2\)
转动起来，则会得到稳定的运动（质点不会跑到无穷远处，而是在一个范围内振动）。求出
\(\omega_2\) 的最小值。

\begin{figure}[H]
\centering
\includegraphics[width=0.88\linewidth,height=\textheight,keepaspectratio,alt={图2-1　旋转抛物面与旋转马鞍面的空间结构}]{../figures/question-02.pdf}
\caption*{图2-1　旋转抛物面与旋转马鞍面的空间结构}
\end{figure}

\subsection{一、题型定位}\label{ux4e00ux9898ux578bux5b9aux4f4d-1}

抛物面底部的两个水平坐标各作简谐运动。换到转动坐标后，离心项和科里奥利项共同改变方程；旋转马鞍面的稳定条件须从完整特征根和临界解判断。

\subsection{二、知识点}\label{ux4e8cux77e5ux8bc6ux70b9-1}

\begin{enumerate}
\def\labelenumi{\arabic{enumi}.}
\tightlist
\item
  小振动近似：在势能极小点附近展开到二阶，等效为简谐运动。
\item
  曲面高度 \(z\simeq\frac a2r^2\) 给出水平方向重力势能
  \(U=\frac12 mga\,r^2\)。
\item
  转动参考系中的惯性力：离心力 \(m\omega^2r\) 与科里奥利力
  \(-2m\omega\,\bm e_z\times\bm v'\)。
\item
  复坐标把两个二阶实方程合并成一个二阶复方程。
\item
  特征方程法解常系数线性微分方程。
\item
  线性有界稳定要求特征根为纯虚数且无产生世俗项的 Jordan
  块；零频临界必须另查。
\end{enumerate}

\subsection{三、物理图像}\label{ux4e09ux7269ux7406ux56feux50cf-1}

在抛物面底部，曲面沿水平方向的高度变化近似为
\(z\simeq\frac a2(x^2+y^2)\)。质点沿 \(x\)
方向微移时重力沿切向的分量正比于位移，指向平衡点，因此 \(x\)、\(y\)
两个方向各自是简谐运动，且对称性使两个方向的频率相同。

换到绕 \(z\)
轴转动的参考系后，质点额外感受到离心力（把质点向外推）与科里奥利力（垂直于相对速度）。科里奥利力不做功，只改变运动方向；离心力与恢复力共同决定径向本征频率。最后用复坐标解出运动，会发现它等于惯性系中的简谐运动再整体叠加一个旋转，频率不变。

马鞍面 \(z=\frac a2x^2-\frac a4y^2\) 的势能为

\[
U=\frac12mga\,x^2-\frac14mga\,y^2.
\]

沿 \(x\) 方向的稳定曲率为 \(ag\)，沿 \(y\) 方向的不稳定曲率大小只有
\(ag/2\)。旋转参考系中的离心力单独看是向外的，并不能被理解成一个``向内的势能井''；真正的稳定化来自离心项与科里奥利速度耦合共同改变特征根。因而必须保留两个方程中的
\(\pm2\omega_2\dot q\)，再由完整特征方程判断稳定性。

\subsection{四、解题入口}\label{ux56dbux89e3ux9898ux5165ux53e3-1}

第（1）问对抛物面写 \(U=\frac12mga(x^2+y^2)\)，由 \(m\ddot x=-mga x\)
读出 \(\omega_0^2=ga\)。

第（2）问分别写出转动系中的三类力：恢复力为 \(-mga\,\bm r\)，离心力为
\(m\omega_1^2\bm r\)，科里奥利力为
\(2m\omega_1(\dot y',-\dot x')\)。将三者分量相加，得到两个耦合方程。

第（3）问令 \(w=X'+iY'\)
把两个方程合并，解特征方程，把解写回惯性系比较。

第（4）问对马鞍面重复第（2）问步骤，写出 \(x'\)、\(y'\) 方程，设解
\(\propto e^{\lambda t}\)，令 \(\lambda\) 为纯虚（或频率平方非负）求得
\(\omega_2\) 的下限。

\subsection{五、解析}\label{ux4e94ux89e3ux6790-1}

\begin{figure}[H]
\centering
\includegraphics[width=0.9\linewidth,height=\textheight,keepaspectratio,alt={图2-2　抛物面小振动与马鞍面旋转稳定化的力平衡}]{../figures/solution-02.pdf}
\caption*{图2-2　抛物面小振动与马鞍面旋转稳定化的力平衡}
\end{figure}

图2-2把离心项和科里奥利项分开，稳定条件由完整特征根决定。等号处另检验重根解，避免把有线性漂移的状态画成稳定振动。

图 2-2 左侧画出抛物面底部 \(x\)、\(y\)
两个方向的恢复力，右侧保留马鞍面上重力曲率、离心项和科里奥利速度耦合，不能把稳定化解释为两个位置力直接抵消。下文
\(\omega_0,\omega_1,\omega_2\) 与图中标注一致。

\needspace{3\baselineskip}
\textbf{（一）第（1）问：抛物面的两个独立简谐频率}

质点在曲面上的势能近似为

\[
U(x,y)=mgz\simeq\frac12mga(x^2+y^2). \tag{1}
\]

沿 \(x\) 方向的运动方程为

\[
m\ddot x=-\frac{\partial U}{\partial x}=-mga\,x. \tag{2}
\]

因此

\[
\ddot x=-ga\,x,
\]

同理 \(\ddot y=-ga\,y\)。两个方向独立、频率相同：

\[
\boxed{\omega_0^2=ga}. \tag{3}
\]

\needspace{3\baselineskip}
\textbf{（二）第（2）问：匀速转动系中的方程}

设转动系绕 \(z\) 轴以角速度 \(\omega_1\) 旋转，\(X',Y'\)
为转动系中的水平坐标。质点在转动系中的绝对加速度等于相对加速度加牵连加速度加科里奥利加速度。等效地，在转动系中列牛顿方程时需加入惯性力：

\begin{itemize}
\tightlist
\item
  离心力 \(m\omega_1^2(X',Y')\)；
\item
  科里奥利力 \(2m\omega_1(\dot Y',-\dot X')\)（因为
  \(\bm\omega_1=\omega_1\bm e_z\)）。
\end{itemize}

恢复力仍为 \(-mga(X',Y')\)。于是

\[
m\ddot X'=-mga\,X'+m\omega_1^2X'+2m\omega_1\dot Y', \tag{4}
\]

\[
m\ddot Y'=-mga\,Y'+m\omega_1^2Y'-2m\omega_1\dot X'. \tag{5}
\]

约去 \(m\)：

\[
\boxed{
\ddot X'=-(ga-\omega_1^2)X'+2\omega_1\dot Y',
}
\]

\[
\boxed{
\ddot Y'=-(ga-\omega_1^2)Y'-2\omega_1\dot X'.
} \tag{6}
\]

\needspace{3\baselineskip}
\textbf{（三）第（3）问：解方程并与（1）对应}

令

\[
w=X'+iY'. \tag{7}
\]

式（6）的第一式加 \(i\) 倍第二式：

\[
\ddot X'+i\ddot Y'
=-(ga-\omega_1^2)(X'+iY')+2\omega_1(\dot Y'-i\dot X').
\]

注意到 \(\dot Y'-i\dot X'=-i(\dot X'+i\dot Y')=-i\dot w\)，于是

\[
\ddot w+2i\omega_1\dot w+(ga-\omega_1^2)w=0. \tag{8}
\]

设 \(w=e^{\lambda t}\)，特征方程为

\[
\lambda^2+2i\omega_1\lambda+(ga-\omega_1^2)=0,
\]

解得

\[
\lambda=-i\omega_1\pm\sqrt{-ga}
=-i\omega_1\pm i\sqrt{ga}. \tag{9}
\]

因此转动系中的通解为

\[
w(t)=e^{-i\omega_1t}\left(Ae^{i\sqrt{ga}\,t}+Be^{-i\sqrt{ga}\,t}\right). \tag{10}
\]

回到惯性系时，把转动系坐标再乘
\(e^{i\omega_1t}\)（逆旋转），得到惯性系中的运动为

\[
w_{\rm lab}(t)=e^{i\omega_1t}w(t)
=Ae^{i\sqrt{ga}\,t}+Be^{-i\sqrt{ga}\,t}. \tag{11}
\]

这正是频率为 \(\sqrt{ga}=\omega_0\)
的简谐运动。结论：\textbf{匀速转动参考系不改变本征频率}，只相当于在解上叠乘一个整体旋转因子
\(e^{i\omega_1t}\)，这与第（1）问的 \(\omega_0\)
完全对应。注意式（10）中频率 \(\sqrt{ga}\) 前的 \(\pm\)
对应两个方向的圆偏振简正模式，它们的频率在转动系中表现为
\(\sqrt{ga}\pm\omega_1\) 的拍频，但惯性系观者仍看到单一的 \(\omega_0\)。

\needspace{3\baselineskip}
\textbf{（四）第（4）问：马鞍面的旋转稳定化}

马鞍面

\[
z=\frac a2x^2-\frac a4y^2
\]

给出势能

\[
U=\frac12mga\,x^2-\frac14mga\,y^2. \tag{12}
\]

沿 \(x\) 方向是稳定简谐（频率平方 \(ga\)），沿 \(y\)
方向是不稳定（``频率平方'' \(-ga/2\)）。在绕 \(z\) 轴以角速度
\(\omega_2\) 转动的参考系中，加入离心力 \(m\omega_2^2(x',y')\)
与科里奥利力 \(2m\omega_2(\dot y',-\dot x')\)，运动方程为

\[
\ddot x'=-(ga-\omega_2^2)x'+2\omega_2\dot y', \tag{13}
\]

\[
\ddot y'=\left(\frac{ga}{2}+\omega_2^2\right)y'-2\omega_2\dot x'. \tag{14}
\]

注意 \(y'\) 方程中位移项符号为``+''，因为 \(y\)
方向势能开口向下；同时其重力项系数是 \(ga/2\)，不能误写成 \(ga\)。

设解为 \((x',y')=(X,Y)e^{\lambda t}\)。代入后得线性齐次方程组：

\[
(\lambda^2+ga-\omega_2^2)X-2\omega_2\lambda Y=0, \tag{15}
\]

\[
2\omega_2\lambda X+
\left(\lambda^2-\frac{ga}{2}-\omega_2^2\right)Y=0. \tag{16}
\]

有非零解要求行列式为零：

\[
(\lambda^2+ga-\omega_2^2)
\left(\lambda^2-\frac{ga}{2}-\omega_2^2\right)
+4\omega_2^2\lambda^2=0.
\]

令 \(q=\lambda^2\)，并记 \(A=ga\)、\(W=\omega_2^2\)。逐项展开：

\[
\begin{aligned}
0
&=(q+A-W)\left(q-\frac A2-W\right)+4Wq\\
&=q^2+\left(\frac A2+2W\right)q
+W^2-\frac{AW}{2}-\frac{A^2}{2}\\
&=q^2+\left(\frac A2+2W\right)q
+(W-A)\left(W+\frac A2\right).
\end{aligned}
\tag{17}
\]

这个关于 \(q\) 的二次方程比直接求四个 \(\lambda\)
更容易判断稳定性。两根之和与积分别为

\[
\begin{aligned}
q_1+q_2&=-\left(\frac A2+2W\right)<0,\\
q_1q_2&=(W-A)\left(W+\frac A2\right). 
\end{aligned}
\tag{18}
\]

判别式为

\[
\begin{aligned}
\Delta_q
&=\left(\frac A2+2W\right)^2
-4(W-A)\left(W+\frac A2\right)\\
&=\frac94A^2+4AW>0. 
\end{aligned}
\tag{19}
\]

所以 \(q_1,q_2\) 总是实数。要求四个 \(\lambda\) 都为纯虚数，就要两个
\(q\) 都为负；由于根和已经为负，只需再要求根积为正：

\[
(W-A)\left(W+\frac A2\right)>0. \tag{20}
\]

又因 \(W+A/2>0\)，故 \(\omega_2^2>ga\)。临界角速度的下界仍为

\[
\boxed{\omega_{2,c}=\sqrt{ga}=\omega_0}. \tag{21}
\]

但题干要求任意给定初态下质点始终在有限范围内振动，临界点必须单独检查。令
\(W=A\)，式（17）化为

\[
q\left(q+\frac52A\right)=0.
\]

\(q=0\)
对应零频模，一般初态会含线性随时间增长的世俗项，不能保证有界。因此严格条件是

\[
\boxed{\omega_2>\sqrt{ga}}. \tag{22}
\]

这里的 \(\sqrt{ga}\)
是稳定区间的下确界，通常称为最小临界角速度。稳定机制是旋转造成的动态耦合，而不是把离心力误当成静态回复力。

\needspace{3\baselineskip}
\textbf{特征根与坐标变换的逐项复核}

第（3）问容易在``本征频率不变''这句话上产生歧义。式（10）在转动坐标中包含两个指数因子

\[
e^{-i(\omega_1-\omega_0)t},\qquad
e^{-i(\omega_1+\omega_0)t},
\]

所以转动系坐标分量实际出现的角频率是 \(|\omega_0-\omega_1|\) 与
\(\omega_0+\omega_1\)。这并不表示物理系统多出了新的弹性常量，而是同一惯性系运动又乘上了坐标轴反向转动因子
\(e^{-i\omega_1t}\)。乘回 \(e^{i\omega_1t}\) 后，两支都恢复为
\(\pm\omega_0\)。因此严谨表述应是：\textbf{曲面的物理小振动频率仍为
\(\omega_0\)，但它在旋转坐标中的分量频率发生了坐标频移。}

还可直接从能量复核科里奥利项的符号。将式（6）第一式乘
\(\dot X'\)、第二式乘 \(\dot Y'\) 后相加，两个速度耦合项为

\[
2\omega_1\dot X'\dot Y'-2\omega_1\dot Y'\dot X'=0,
\]

说明科里奥利力瞬时功率恒为零。若两式的科里奥利项写成同号，这一抵消立即失败，方程必错。离心项则可并入转动系有效势

\[
U_{\rm eff}=U-\frac12m\omega_1^2(X'^2+Y'^2),
\]

注意这里是减号，正好说明``离心势能把曲率翻正''的说法不成立。

对第（4）问，令 \(\omega_2^2=ga+\varepsilon\)，其中 \(\varepsilon>0\)
很小。式（17）的常数项为

\[
(W-A)\left(W+\frac A2\right)
=\varepsilon\left(\frac32A+\varepsilon\right)>0,
\]

而根和为负，所以两根均为负。随着 \(\varepsilon\to0^+\)，其中一个 \(q\)
趋于零，相应慢模频率趋于零，这正是临界点出现世俗漂移的前兆；当
\(\varepsilon<0\) 时根积变负，必有一个
\(q>0\)，于是出现实指数增长。这个极限检验同时核对了临界值、严格不等号和稳定性的物理含义。

\needspace{3\baselineskip}
\textbf{临界转速并不产生一般的有界解}

严格稳定条件是
\(\omega_2>\sqrt{ga}\)，所以题问的``最小值''应理解为下确界
\(\sqrt{ga}\)。在等号处令 \(\Omega=\sqrt{ga}\)，转动系方程允许

\[
x'=vt+x_0,\qquad y'=\frac{4v}{3\Omega},
\]

其中 \(v\) 任意。代回两方程即可核对，然而 \(x'\)
随时间增长。因此只检查特征频率平方非负会错误地把临界也归入稳定区。曲面线性化还要求振幅远小于
\(1/a\)，结论是底部附近的线性稳定，不是任意振幅的全局稳定保证。

\subsection{六、关键洞察}\label{ux516dux5173ux952eux6d1eux5bdf-1}

解析图把离心项和科里奥利项分开，稳定条件由完整特征根决定。等号处另检验重根解，避免把有线性漂移的状态画成稳定振动。

第（4）问的妙处在于动态稳定：两个方向的曲率并不对称，稳定曲率为
\(ga\)，不稳定曲率大小为
\(ga/2\)；但离心项与科里奥利耦合后，根积恰好含因子
\((\omega_2^2-ga)(\omega_2^2+ga/2)\)。稳定阈值必须从完整四阶特征式得到，不能只比较某一坐标方向的静态``势能曲率''。

临界值附近的一快一慢两支模式是本题最值得保留的物理图像。\(\omega_2\)
刚超过 \(\sqrt{ga}\)
时，慢模负责大尺度缓慢摆动，快模叠加细小快速振动；若恰在临界点，慢模退化成零频并可能出现线性漂移。因此竞赛答案写``最小值
\(\sqrt{ga}\)''时，教师讲解必须补一句：它是稳定区间的临界下界，要求任意初态都在有限范围内振动时应取严格大于。

\subsection{七、易错点}\label{ux4e03ux6613ux9519ux70b9-1}

\begin{enumerate}
\def\labelenumi{\arabic{enumi}.}
\tightlist
\item
  把抛物面高度系数写错。\(z=\frac a2(x^2+y^2)\) 时势能系数是
  \(\frac12mga\)，频率平方是 \(ga\)，不是 \(2ga\)。
\item
  转动系方程中漏掉科里奥利力，或把 \(\dot Y'\)
  项的符号写反。可检验：科里奥利力不做功，两个方程必须使
  \(\frac d{dt}(X'^2+Y'^2)\) 中不含耦合项贡献。
\item
  特征方程解错符号，导致把 \(ga-\omega_1^2\) 误当频率平方。用
  \(\omega_1=0\) 极限核对：式（8）应退化为 \(\ddot w+ga\,w=0\)。
\item
  第（4）问把马鞍面误抄成对称形式 \(a(x^2-y^2)/2\)，从而把 \(y\)
  方向的不稳定曲率写成 \(ga\)。题设实际给出 \(-ay^2/4\)，故该曲率大小为
  \(ga/2\)。
\item
  稳定条件误用为``频率为正''，正确判据是``频率平方非负''（\(\lambda\)
  纯虚）。
\end{enumerate}

验算清单依次为：\(\omega_1=0\)
时退回普通小振动；两条科里奥利项做功相消；马鞍面 \(y\) 方程在
\(\omega_2=0\) 时退化为
\(\ddot y=(ga/2)y\)；\(\omega_2^2=ga+\varepsilon\) 时两个
\(q=\lambda^2\) 均为负。临界等号还要检查零根的重数，不能仅看到``没有正的
\(q\)''就断言一般运动有界。

\subsection{八、课堂讲解建议}\label{ux516bux8bfeux5802ux8bb2ux89e3ux5efaux8bae-1}

先让学生从势能 \(U=\frac12mga(x^2+y^2)\)
独立推出式（3），再演示复坐标代换。第（4）问应把离心项与科里奥利项同时保留，避免用静态势能类比替代特征根分析。最后用数值例子检验：取
\(ga=10\ \mathrm{s^{-2}}\)，分别令 \(\omega_2=3\ \mathrm{s^{-1}}\) 与
\(4\ \mathrm{s^{-1}}\)，比较特征值是实数还是纯虚数。

板书可按``几何---势能---方程---特征根''四层推进。第一层只解释曲面高度为什么能近似成二次型；第二层让学生从
\(U\)
的曲率读出稳定、失稳方向；第三层再加入转动系惯性力；第四层才进入行列式。这样学生会明白特征方程不是凭空出现，而是在量化势能曲率与速度耦合的竞争。

一个有效的即时练习是把两条科里奥利项中的一个符号故意写反，让学生用``科里奥利力不做功''检查。随后把
\(\omega_2^2=ga+\varepsilon\)
代入，让学生读出快、慢两支频率。慢频率趋零的图像比只背
\(\omega_2>\sqrt{ga}\) 更能解释临界点为何不能保证一般初态有界。

\subsection{九、变式训练}\label{ux4e5dux53d8ux5f0fux8badux7ec3-1}

\begin{enumerate}
\def\labelenumi{\arabic{enumi}.}
\tightlist
\item
  若抛物面改为 \(z=\frac b2(x^2+y^2)\)，重求 \(\omega_0\) 与
  \(\omega_{2,\min}\)。
\item
  第（2）问中若初始相对速度不为零，解出通解并给出 \(X',Y'\)
  的显式时间依赖。
\item
  讨论一般曲面
  \(z=\frac12(\alpha x^2+\beta y^2)\)（\(\alpha\beta<0\)）旋转稳定化的条件。
\end{enumerate}

\textbf{变式解析要点：} 第 1 项只需把 \(a\) 换成 \(b\)，得到
\(\omega_0=\sqrt{gb}\)。第 2 项仍用

\[
w(t)=e^{-i\omega_1t}\left(Ae^{i\omega_0t}+Be^{-i\omega_0t}\right),
\]

常数 \(A,B\) 由 \(w(0)\)、\(\dot w(0)\) 联立确定。第 3
项不应试图把各向异性二次型直接写成复数单方程；应设
\((x,y)=(X,Y)e^{\lambda t}\)，保留矩阵

\[
\begin{pmatrix}
\lambda^2+g\alpha-\Omega^2&-2\Omega\lambda\\
2\Omega\lambda&\lambda^2+g\beta-\Omega^2
\end{pmatrix},
\]

令行列式为零，再要求所有 \(\lambda\) 的实部为零。取本题
\(\alpha=a,\beta=-a/2\) 时，常数项应分解为

\[
(\Omega^2-ga)\left(\Omega^2+\frac{ga}{2}\right),
\]

这是一般式的必要复核。

\newpage

\section{第 3 题}\label{ux7b2c-3-ux9898}

如图，质量 \(m\)、长度 \(2l\)
的匀质薄板在光滑桌面上，两桌面之间有两个以角速度 \(\omega\)
匀速转动的半径为 \(r\) 的圆柱，摩擦系数都为
\(\mu\)，桌面边缘和两圆柱上顶点比桌面稍高一点点，使得薄板总是点接触。

（1）若将匀质薄板放在两圆柱上，并使得薄板中心稍微偏离两圆柱中心，求证当
\(\omega>0\) 时薄板将做简谐运动，并求其周期；

（2）当 \(\omega=0\) 时，求匀质板右侧能顺利到达右侧桌面的最小初始速度
\(v_{10}\)；

（3）接（1），求这种情况下匀质板右侧到达左桌面边缘，到匀质板右侧到达右桌面边缘的时间；

（4）此问不是考题不计分。\(\omega>0\) 且初始速度为
\(v_0=2\omega r\)，令无量纲参量
\(\beta=\dfrac{\mu gl}{\omega^2r^2}\)，在 \(\beta\)
取各种可能值的情况下，求匀质板右侧到达左桌面边缘，到匀质板左侧到达右桌面边缘的时间。

\begin{figure}[H]
\centering
\includegraphics[width=0.66\linewidth,height=\textheight,keepaspectratio,alt={图3-1　薄板搭在两转动圆柱上}]{../figures/question-03.pdf}
\caption*{图3-1　薄板搭在两转动圆柱上}
\end{figure}

\subsection{一、题型定位}\label{ux4e00ux9898ux578bux5b9aux4f4d-2}

板心位置决定支持力，板速与圆柱表面速度的大小决定摩擦方向。两类事件都要分段；在相对速度为零时，还可能出现静摩擦维持的恒速区间。

\subsection{二、知识点}\label{ux4e8cux77e5ux8bc6ux70b9-2}

\begin{enumerate}
\def\labelenumi{\arabic{enumi}.}
\tightlist
\item
  动摩擦力方向与相对滑动速度反向，大小 \(\mu N\)。
\item
  力矩平衡求两支撑点法向力随板心位移的变化。
\item
  线性恢复力 \(F=-kx\) 与简谐运动周期 \(T=2\pi\sqrt{m/k}\)。
\item
  恒力做功与动能定理。
\item
  简谐运动的相位积分 \(t=\int dx/\sqrt{...}\) 或
  \(x(t)=A\sin(\omega_0 t+\varphi)\)。
\end{enumerate}

\subsection{三、物理图像}\label{ux4e09ux7269ux7406ux56feux50cf-2}

固定在空间中的四个关键位置依次是左桌边、左圆柱顶点、右圆柱顶点和右桌边，相邻间距均为
\(l\)；两圆柱中心的水平距离是 \(l\)，而薄板长度是
\(2l\)。板心在两圆柱中点右侧 \(x\)
处时，右支撑力增大、左支撑力减小。两顶点的表面速度相向，因而左接触处摩擦向右、右接触处摩擦向左；右侧较大的摩擦力使合力指回平衡位置。

薄板从左桌面向右运动时，支承状态按``无摩擦---左桌边与左圆柱---两圆柱---右圆柱与右桌边---无摩擦''依次切换。第（2）至第（4）问的路程、摩擦功和时间必须按这些区间分别计算。

\subsection{四、解题入口}\label{ux56dbux89e3ux9898ux5165ux53e3-2}

第（1）问先由力矩平衡写 \(N_L,N_R\)，再由两摩擦力之差证明
\(F_x=-2\mu mgx/l\)。第（2）问不能把全程都当作摩擦力 \(\mu mg\)
的匀减速：第一段过渡区的圆柱支持力从 \(0\) 线性增到
\(mg\)，必须先积分求功。第（3）问按无摩擦、单圆柱过渡和两圆柱支承三段求时间。第（4）问则用
\(v=\omega r\) 作为摩擦换向的事件界面，对五个空间区间分段积分。

\subsection{五、解析}\label{ux4e94ux89e3ux6790-2}

\begin{figure}[H]
\centering
\includegraphics[width=0.891\linewidth,height=\textheight,keepaspectratio,alt={图3-2　两圆柱支承时的受力关系、线性回复力以及整板横跨缝隙时的五个运动区间。}]{../figures/solution-03.pdf}
\caption*{图3-2　两圆柱支承时的受力关系、线性回复力以及整板横跨缝隙时的五个运动区间。}
\end{figure}

图3-2的支撑长度给出 \(N_L,N_R\)，摩擦方向还由 \(v-u\)
的符号决定。下面的分段按接触区域和速度事件一起切换。

图 3-2 左上图给出板心偏离 \(x\) 时两支撑力 \(N_L,N_R\)
的分配，右上图把两块动摩擦合成线性净力；下方横轴使用板心坐标
\(X\)，第（2）至第（4）问将直接按这五个区间列式。

\needspace{3\baselineskip}
\textbf{（一）第（1）问：支撑力分配与简谐运动}

设板心相对两圆柱中点向右偏离 \(x\)，两圆柱位于
\(x=\pm l/2\)。左、右支撑点到板心的水平力臂分别为
\(l/2+x\)、\(l/2-x\)。竖直平衡与绕板心取矩：

\[
N_L+N_R=mg, \tag{1}
\]

\[
N_L\left(\frac l2+x\right)=N_R\left(\frac l2-x\right). \tag{2}
\]

解得

\[
N_L=\frac{mg}{l}\left(\frac l2-x\right),\qquad
N_R=\frac{mg}{l}\left(\frac l2+x\right). \tag{3}
\]

两圆柱反向转动，且 \(|v|<\omega r\)
时：左圆柱顶部向右运动，板相对左接触点向左滑，动摩擦对板向右
\(\mu N_L\)；右圆柱顶部向左运动，板相对右接触点向右滑，动摩擦对板向左
\(\mu N_R\)。水平方向：

\[
m\ddot x=\mu N_L-\mu N_R
=-\frac{2\mu mg}{l}x. \tag{4}
\]

这是线性恢复力，角频率

\[
\omega_0=\sqrt{\frac{2\mu g}{l}}, \tag{5}
\]

周期

\[
\boxed{T=2\pi\sqrt{\frac{l}{2\mu g}}}. \tag{6}
\]

\needspace{3\baselineskip}
\textbf{（二）第（2）问：\(\omega=0\) 时的最小初速度}

以两圆柱中点为原点，用 \(X\)
表示板心坐标。左桌边、左圆柱、右圆柱、右桌边分别位于

\[
X=-\frac{3l}{2},\quad-\frac l2,\quad\frac l2,\quad\frac{3l}{2}.
\]

初始时板右端在左桌边，板心在 \(X=-5l/2\)。板先在光滑左桌面上前进
\(l\)，这一段无摩擦功。

当 \(X\) 从 \(-3l/2\) 增至 \(-l/2\)
时，板由左桌边与左圆柱共同支承。记左圆柱的支持力为
\(N_1\)。竖直力平衡及力矩平衡给出

\[
N_1=mg\frac{X+3l/2}{l}.
\]

因此这一过渡段需克服的摩擦功为

\[
\begin{aligned}
W_1
&=\int_{-3l/2}^{-l/2}\mu N_1\,dX\\
&=\frac{\mu mg}{l}\int_{-3l/2}^{-l/2}
\left(X+\frac{3l}{2}\right)dX\\
&=\frac{\mu mg}{l}
\left[\frac12\left(X+\frac{3l}{2}\right)^2\right]_{-3l/2}^{-l/2}\\
&=\frac12\mu mgl. 
\end{aligned}
\tag{7}
\]

接着，板心从 \(X=-l/2\) 运动到
\(X=l/2\)，两圆柱同时支承薄板。两支持力之和恒为
\(mg\)，因而摩擦力总大小为 \(\mu mg\)，这段距离为 \(l\)：

\[
W_2=\mu mg\,l. \tag{8}
\]

``刚好到达''的临界情形中，终速度为零。初动能等于两段摩擦功：

\[
\begin{aligned}
\frac12mv_{10}^{2}
&=W_1+W_2\\
&=\frac12\mu mgl+\mu mgl\\
&=\frac32\mu mgl.
\end{aligned}
\]

约去 \(m/2\) 后得

\[
\boxed{v_{10}=\sqrt{3\mu gl}}. \tag{9}
\]

\needspace{3\baselineskip}
\textbf{（三）第（3）问：一般初速的事件积分}

本问的终点为板右端到达右桌边，对应板心从 \(X=-5l/2\) 到 \(X=l/2\)。令
\(u=\omega r>0\)，在向右运动的区间，把 \(w=v^2\) 作为位置函数，有
\(dw/dX=2a(X,v)\)。

在 \(-5l/2<X<-3l/2\)，无摩擦，\(w=v_0^2\)。在左过渡区 \(-3l/2<X<-l/2\)，

\[
a=-\frac{\mu g}{l}(X+3l/2)\operatorname{sgn}(v-u).
\]

在两圆柱区 \(-l/2<X<l/2\)，

\[
a=\begin{cases}-\mu g,&v>u,\\-2\mu gX/l,&0<v<u.\end{cases}
\]

若 \(v=u\) 且
\(X\le0\)，左接触的静摩擦可以平衡右接触的滑动摩擦，板可保持 \(v=u\)；到
\(X=0\) 后静摩擦不足，才进入 \(v<u\)
的简谐减速段。这里按静摩擦上限不小于给定滑动摩擦处理，若另给不同静摩擦因数，应重新核对恒速区间。

从 \(w(-5l/2)=v_0^2\) 积分，每遇支撑变化、\(w=u^2\) 或 \(w=0\)
都重新选分支。到达终点且沿途 \(w>0\) 时，所求时间是

\[
\boxed{t=\int_{-5l/2}^{l/2}\frac{dX}{\sqrt{w(X)}}}.
\]

各段 \(w\)
为常数、一次或二次多项式，积分分别为匀速、反三角或反双曲函数。若轨迹在终点前折返，第一次向右穿越不存在，应继续按反向滑动的方程判断，而不能对负的
\(w\) 开方。

例如 \(v_0^2\ge u^2+3\mu gl\) 时，全程保持 \(v\ge u\)，左过渡区末速
\(v_1=\sqrt{v_0^2-\mu gl}\)，双圆柱区为匀减速。此时可直接写

\[
\boxed{t=\frac l{v_0}+\sqrt{\frac l{\mu g}}\arcsin\frac{\sqrt{\mu gl}}{v_0}
+\frac{\sqrt{v_0^2-\mu gl}-\sqrt{v_0^2-3\mu gl}}{\mu g}}.
\]

原整理稿把第二段的 \(v>u\) 减速解与第三段的 \(v<u\)
简谐解直接相接，未检查换向事件，所得单一公式不能用于一般
\(v_0\)。上面的事件积分给出一般解的确定步骤。

\needspace{3\baselineskip}
\textbf{（四）第（4）问：\(v_0=2\omega r\) 时的完整分段结果}

记 \(u=\omega r\)，则 \(v_0=2u\)，且

\[
\beta=\frac{\mu gl}{u^2},\qquad \beta>0. \tag{17}
\]

本问的起点是板右端在左桌边，即 \(X=-5l/2\)；终点是板左端在右桌边，即
\(X=5l/2\)。图 3-2 下方的五个区间恰好是必须累加的五段时间。

第 1 段无摩擦，速度为 \(2u\)：

\[
t_1=\frac{l}{2u}. \tag{18}
\]

先考虑 \(\beta<3\)。第 2 段仍用 \(y=X+3l/2\)，由于 \(\beta<3\)，到达
\(y=l\) 时的速度

\[
v_1=u\sqrt{4-\beta}>u,
\]

因此这一整段中薄板始终比左圆柱顶面更快，摩擦向左，式（11）可用。令
\(A_1=2l/\sqrt\beta\)，得

\[
t_2=\frac{l}{u\sqrt\beta}
\arcsin\frac{\sqrt\beta}{2}. \tag{19}
\]

第 3 段刚进入两圆柱区时 \(v_1>u\)，两个接触点的摩擦都向左，所以暂时有
\(\ddot X=-\mu g\)。从 \(X=-l/2\) 量起位移 \(s\)，

\[
\frac{v^2}{u^2}=4-\beta-2\beta\frac{s}{l}. \tag{20}
\]

若 \(0<\beta<1\)，在 \(s=l\) 时仍有

\[
v_2=u\sqrt{4-3\beta}>u,
\]

因而第 3 段全程为匀减速，

\[
t_3=\frac{v_1-v_2}{\mu g}
=\frac{l}{\beta u}
\left(\sqrt{4-\beta}-\sqrt{4-3\beta}\right). \tag{21}
\]

第 4 段令 \(z=X-3l/2\)，则 \(z\) 从 \(-l\) 增至 \(0\)。右圆柱的支持力为
\(mg(-z/l)\)，摩擦向左，故

\[
\ddot z=\frac{\mu g}{l}z=\omega_1^2z. \tag{22}
\]

以进入该段的时刻为新的 \(t=0\)，\(z(0)=-l, \dot z(0)=v_2\)，解为

\[
z=-l\cosh\omega_1t
+\frac{v_2}{\omega_1}\sinh\omega_1t. \tag{23}
\]

令 \(z=0\)，得

\[
t_4=\frac1{\omega_1}
\operatorname{artanh}\frac{\omega_1l}{v_2}
=\frac{l}{u\sqrt\beta}
\operatorname{artanh}\frac{\sqrt\beta}{\sqrt{4-3\beta}}. \tag{24}
\]

根号内比较可见，只有 \(\beta<1\) 时
\(v_2>\omega_1l\)，薄板才能穿过这一段。由式（23）或动能定理，到达
\(z=0\) 时的速度为

\[
v_3=\sqrt{v_2^2-\omega_1^2l^2}
=2u\sqrt{1-\beta}. \tag{25}
\]

第 5 段薄板已离开圆柱，再以恒速 \(v_3\) 前进 \(l\)：

\[
t_5=\frac{l}{2u\sqrt{1-\beta}}. \tag{26}
\]

因此，当 \(0<\beta<1\) 时，

\[
\boxed{
\begin{aligned}
t=\frac{l}{u}\Bigg[&\frac12
+\frac1{\sqrt\beta}\arcsin\frac{\sqrt\beta}{2}\\
&+\frac{\sqrt{4-\beta}-\sqrt{4-3\beta}}{\beta}\\
&+\frac1{\sqrt\beta}\operatorname{artanh}
\frac{\sqrt\beta}{\sqrt{4-3\beta}}
+\frac1{2\sqrt{1-\beta}}\Bigg].
\end{aligned}} \tag{27}
\]

当 \(\beta=1\)
时，薄板到达右桌边附近时速度趋于零，不能再走完最后的光滑段，式（27）也显示时间发散。

当 \(1<\beta<3\) 时，双圆柱区内首次达到 \(v=u\) 的位置为

\[
X_*=\frac{3-2\beta}{2\beta}l.
\]

必须按 \(X_*\) 的正负再分段。

\begin{itemize}
\tightlist
\item
  \(1<\beta<3/2\)：\(X_*>0\)，立即进入 \(v<u\)
  的简谐减速，右圆柱位置的速度满足
  \(v_R^2/u^2=(3\beta^2-10\beta+9)/(2\beta)\)。右过渡区所需能量对应
  \(\mu gl=\beta u^2\)，而
  \(v_R^2-\beta u^2=u^2(\beta-1)(\beta-9)/(2\beta)<0\)，不能穿过。
\item
  \(3/2\le\beta<2\)：\(X_*\le0\)，先由静摩擦维持 \(v=u\) 到
  \(X=0\)，再作简谐减速。于是
  \(v_R^2/u^2=1-\beta/2\)，仍不足以通过右过渡区。
\item
  \(\beta\ge2\)：若 \(\beta<3\)，先在双圆柱区达到 \(v=u\)；若
  \(\beta\ge3\)，更早在左过渡区达到 \(v=u\)，位置为
  \(X+3l/2=l\sqrt{3/\beta}\)。均可保持恒速到 \(X=0\)，此后转向位置为
  \(X_{\rm turn}=l/\sqrt{2\beta}\le l/2\)，连右过渡区都不能完整进入。
\end{itemize}

在
\(\beta=1\)，右过渡区末端只在无限时间极限达到，最后光滑段无法走完。因此

\[
\boxed{0<\beta<1\text{ 时式（27）给出有限通过时间；}\quad\beta\ge1\text{ 时不能完成穿越。}}
\]

折返后圆柱虽能向板输入能量，却不会把进入右过渡区的右行速度提高到 \(u\)
以上；\(\beta>1\) 时通过该区需要速度大于
\(\sqrt\beta\,u>u\)，故以后也不能完成穿越。

\needspace{3\baselineskip}
\textbf{（五）结果复核}

第一，将式（3）相加得 \(N_L+N_R=mg\)，相减得
\(N_R-N_L=2mgx/l\)；竖直平衡和合力斜率同时被满足。周期的量纲为

\[
\left[\sqrt{\frac{l}{g}}\right]=\mathrm s,
\]

且质量 \(m\) 消去，与动摩擦和惯性都正比于 \(m\) 一致。

第二，第（2）问的总摩擦功是

\[
W_1+W_2=\frac32\mu mgl.
\]

若误把全部有接触的 \(2l\) 都乘以 \(\mu mg\)，会得到
\(2\mu mgl\)，这与过渡段支持力从零递增的事实冲突。

第三，式（27）在 \(\beta\to0\) 时的五项依次趋于
\(1/2,1/2,1/2,1/2,1/2\)，因而

\[
\lim_{\beta\to0}t
=\frac{5l}{2u}.
\]

这正是无摩擦极限下以恒速 \(2u\) 走完总路程 \(5l\)
所需的时间，说明五段路程没有遗漏。当 \(\beta\to1^-\)
时，式（24）和式（26）同时发散，与薄板在临界状态下无法完成穿越的物理图像一致。

\needspace{3\baselineskip}
\textbf{简谐周期的条件}

第一问 \(T=2\pi\sqrt{l/(2\mu g)}\)
仅在板一直搭在两圆柱上且两摩擦方向相反时成立。若振幅为 \(A\)，应有
\(A<l/2\) 和
\(A\sqrt{2\mu g/l}<u\)。图3-2把支持区间和速度分支分开标出；恒速平台是静摩擦阶段，不能画成瞬间反向的动摩擦。

\subsection{六、关键洞察}\label{ux516dux5173ux952eux6d1eux5bdf-2}

解析图的支撑长度给出 \(N_L,N_R\)，摩擦方向还由 \(v-u\)
的符号决定。后续的分段按接触区域和速度事件一起切换。

\[
k_{\mathrm{eff}}=\frac{2\mu mg}{l}.
\]

简谐运动并不是来自材料的弹性形变，而是来自两支撑力随板心位置线性重新分配。学生若只记住周期而没有先得到
\(N_R-N_L=2mgx/l\)，就没有真正掌握模型。

第（2）问的关键不是``总路程乘恒摩擦力''，而是过渡段中圆柱支持力从 \(0\)
增到
\(mg\)，所以摩擦功必须积分。第（3）、（4）问的关键则是把空间接触区间与速度事件同时作为分段边界：跨过
\(X=-3l/2,-l/2,l/2,3l/2\) 会改变支承状态，跨过 \(v=\omega r\)
会改变摩擦方向。

\subsection{七、易错点}\label{ux4e03ux6613ux9519ux70b9-2}

\begin{enumerate}
\def\labelenumi{\arabic{enumi}.}
\tightlist
\item
  把``板长 \(2l\)''误当成``圆柱间距 \(2l\)''。圆柱间距实际为
  \(l\)，这会使回复力斜率和周期中的因子 \(2\) 全部改变。
\item
  把两力臂写成 \(l\pm x\)。正确力臂是 \(l/2\pm x\)。
\item
  在第（2）问中把所有接触路段都当成 \(\mu mg\)
  的恒摩擦，漏掉桌边---圆柱过渡段的线性支持力。
\item
  在第（3）问直接套半周期，既没有核对起点、终点，也没有比较板速与
  \(\omega r\)。
\item
  在第（4）问只按位置分段而不按 \(v=\omega r\)
  分段，或只按速度换向而漏掉支承点切换。
\item
  把 \(\operatorname{artanh}\) 误写成
  \(\arctan\)。右过渡段满足指数型方程
  \(\ddot z=\omega_1^2z\)，不是简谐方程。
\item
  得到 \(\beta\ge1\) 时仍强行给出有限时间；实际上薄板不能到达规定终点。
\end{enumerate}

批改时先检查学生是否在图 3-2
的横轴上标出六个板心位置，再检查每段有效支持点和摩擦方向。第（2）问必须出现
\(W_1=\mu mgl/2\)、\(W_2=\mu mgl\)；第（4）问必须明确写出 \(0<\beta<1\)
才能通过，并用 \(\beta\to0\) 的极限核对总路程 \(5l\)。

\subsection{八、课堂讲解建议}\label{ux516bux8bfeux5802ux8bb2ux89e3ux5efaux8bae-2}

第一轮讲解只处理板同时搭在两圆柱上的状态。先把板心放在中点右侧，让学生判断右支持力较大，再沿图中两顶面速度判断两摩擦方向，最后才列力矩方程。板书应保留

\[
N_R-N_L=\frac{2mg}{l}x
\quad\Longrightarrow\quad
F_x=-\frac{2\mu mg}{l}x
\]

这一条主线，并让学生说明负号为何表示回复力。

第二轮讲解沿图 3-2
下方的横轴移动一个``板心标记''。每跨过一个刻度，就让学生重新说出当前有哪些支撑点；这样第（2）问的两个摩擦功不会被混为一段。积分
\(W_1\) 时可先画 \(N_1\)---\(X\) 直线，用三角形面积 \(\frac12mgl\)
作一次快速复核。

第三轮讲解第（4）问时，将``位置事件''和``速度事件''用两种颜色记录。先完整推导
\(0<\beta<1\) 的五段时间，再处理 \(1<\beta<3\) 中的换向位置
\(X_*\)，最后用能量不足证明不能穿过右过渡段。这样学生看到的是一套可复用的分段动力学方法，而不是一串孤立公式。

\subsection{九、变式训练}\label{ux4e5dux53d8ux5f0fux8badux7ec3-2}

\begin{enumerate}
\def\labelenumi{\arabic{enumi}.}
\tightlist
\item
  若两圆柱同向转动，重新分析净摩擦并判断能否形成简谐运动。
\item
  若板心初始偏离为 \(x_0=l/2\)，求振幅和到达另一圆柱上方的时刻。
\item
  保留一般圆柱间距 \(d\)，重写支撑力分配与周期。
\end{enumerate}

\textbf{变式解析要点：}
两圆柱同向转动时，若板速小于两接触点表面速度，两处摩擦同向，合力不再是指向中点的线性回复力，因而通常不能形成正文中的简谐运动。第
2 项若 \(x(0)=l/2\) 且 \(v(0)=0\)，振幅就是
\(l/2\)；到达另一侧对称位置需半周期。第 3 项令两支点位于
\(\pm d/2\)，竖直平衡与力矩平衡给出

\[
N_L=mg\left(\frac12-\frac{x}{d}\right),\qquad
N_R=mg\left(\frac12+\frac{x}{d}\right),
\]

所以

\[
F_x=-\frac{2\mu mg}{d}x,\qquad
T=2\pi\sqrt{\frac{d}{2\mu g}}.
\]

代入本题的圆柱间距 \(d=l\)，即还原正文结果。

\newpage

\section{第 4 题}\label{ux7b2c-4-ux9898}

焦距为 \(f\) 的理想薄凸透镜对一个正三角形 \(ABC\) 成像。正三角形的边
\(AC\) 垂直于主轴，顶点 \(A\) 位于主轴上，顶点 \(B\) 位于 \(AC\)
与透镜之间。要求三个顶点的像 \(A'B'C'\) 构成内角为
\(30^\circ\)、\(60^\circ\)、\(90^\circ\) 的直角三角形。

（1）当像三角形的一条直角边垂直于主轴时，求 \(A\) 点物距的较大值
\(u_1\)；

（2）求（1）所对应的正三角形边长 \(l_1\)；

（3）仍令像三角形的一条直角边垂直于主轴，求 \(A\) 点物距的较小值
\(u_2\)；

（4）求（3）所对应的正三角形边长 \(l_2\)；

（5）当像三角形的斜边垂直于主轴时，求 \(A\) 点物距 \(u_3\)；

（6）求（5）所对应的正三角形边长 \(l_3\)。

\begin{figure}[H]
\centering
\includegraphics[width=0.72\linewidth,height=\textheight,keepaspectratio,alt={图4-1　正三角形各顶点相对薄凸透镜与主轴的位置关系。}]{../figures/question-04.pdf}
\caption*{图4-1　正三角形各顶点相对薄凸透镜与主轴的位置关系。}
\end{figure}

\subsection{一、题型定位}\label{ux4e00ux9898ux578bux5b9aux4f4d-3}

三个顶点的物距不同，必须逐点成像。先求各像点坐标，再用直角和边长比约束像三角形，不能把原三角形按一个倍率缩放。

\subsection{二、知识点}\label{ux4e8cux77e5ux8bc6ux70b9-3}

\begin{enumerate}
\def\labelenumi{\arabic{enumi}.}
\tightlist
\item
  薄透镜成像公式 \[
  \frac1p+\frac1q=\frac1f,
  \qquad
  q=\frac{fp}{p-f}.
  \]
\item
  实物、实像情况下的横向放大率绝对值 \[
  m=\frac qp=\frac{f}{p-f}.
  \]
\item
  正三角形的水平投影与竖直坐标： \[
  s=\frac{\sqrt3}{2}l,
  \qquad
  y_B=\frac l2.
  \]
\item
  直角判据 \(\bm a\cdot\bm b=0\)，以及
  \(30^\circ\)-\(60^\circ\)-\(90^\circ\) 三角形的边长比 \(1:\sqrt3:2\)。
\item
  物理根筛选：所有物点均位于焦点外侧，且图像的顶点顺序与题设相符。
\end{enumerate}

\subsection{三、物理图像}\label{ux4e09ux7269ux7406ux56feux50cf-3}

取透镜光心为原点，主轴为 \(x\) 轴，物方在左侧。令

\[
A=(-u,0),\qquad
C=(-u,l),\qquad
B=(-u+s,l/2),
\qquad
s=\frac{\sqrt3}{2}l.
\]

因为 \(A,C\) 位于同一物平面，它们的像 \(A',C'\) 位于同一像平面，所以
\(A'C'\) 必垂直于主轴。\(B\) 比 \(A,C\) 更靠近透镜，物距为
\(u-s\)，其像距和放大率均与 \(A,C\) 不同。题目两种情形实质上是：

\begin{itemize}
\tightlist
\item
  \(A'C'\) 是直角边；
\item
  \(A'C'\) 是斜边。
\end{itemize}

\subsection{四、解题入口}\label{ux56dbux89e3ux9898ux5165ux53e3-3}

把 \(A'\) 作为像方坐标原点，只保留三个几何量：

\[
h_A=|A'C'|,\qquad
h_B=\text{$B'$ 到主轴的距离},\qquad
d=\text{$B'$ 与 $A'C'$ 的水平距离}.
\]

先由薄透镜公式把 \(h_A,h_B,d\) 写成 \(u,l,f\)，再分别使用``\(B'\) 与
\(C'\)
等高''或``\(A'B'\perp C'B'\)''。这样可把光学问题变成一个清晰的平面几何问题。

\subsection{五、解析}\label{ux4e94ux89e3ux6790-3}

图4-2按三组解绘制像三角形。下面先计算物方投影 \(s\)，再求像方的
\(h_A,h_B,d\)；这些量决定图中的三条边和直角位置。

\begin{figure}[H]
\centering
\includegraphics[width=0.9\linewidth,height=\textheight,keepaspectratio,alt={图4-2　三组按薄透镜坐标绘制的像三角形}]{../figures/solution-04.pdf}
\caption*{图4-2　三组按薄透镜坐标绘制的像三角形}
\end{figure}

图4-2的三组像三角形按逐点薄透镜坐标绘制，同一小图横纵比例一致。前两组的竖直像边都是直角边，分别较长和较短；第三组竖直像边为斜边，直角在B点。

\needspace{3\baselineskip}
\textbf{（一）建立三个像点的坐标}

记

\[
a=u-f.
\]

则 \(A,C\) 的物距均为 \(u\)，像距为

\[
q_A=\frac{fu}{u-f}=f+\frac{f^2}{a}. \tag{1}
\]

\(A\) 在主轴上，所以可取

\[
A'=(q_A,0).
\]

\(C\) 的物高为 \(l\)，故其倒立像位于主轴下方，像高绝对值为

\[
h_A=\frac{fl}{u-f}=\frac{fl}{a}. \tag{2}
\]

于是

\[
C'=(q_A,-h_A).
\]

点 \(B\) 比 \(A\) 靠近透镜的水平距离为

\[
s=\frac{\sqrt3}{2}l,
\]

所以 \(B\) 的物距为 \(u-s\)。为保证 \(B\) 也成实像，需有 \(u-s>f\)，即
\(a>s\)。其像距为

\[
q_B=\frac{f(u-s)}{u-s-f}
=f+\frac{f^2}{a-s}. \tag{3}
\]

\(B\) 的物高为 \(l/2\)，所以

\[
h_B=\frac{fl}{2(a-s)}. \tag{4}
\]

因此

\[
B'=(q_B,-h_B).
\]

\(B'\) 到直线 \(A'C'\) 的水平距离为

\[
\begin{aligned}
d
&=q_B-q_A\\
&=\frac{f^2}{a-s}-\frac{f^2}{a}\\
&=\frac{f^2s}{a(a-s)}.
\end{aligned}
\tag{5}
\]

式（2）、（4）、（5）是全题的共同基础。

\needspace{3\baselineskip}
\textbf{（二）\(A'C'\) 为直角边时的几何条件}

\(A'C'\) 已经垂直主轴。若它是直角边，则另一条直角边必须水平。\(B'\)
不可能落在主轴上，因此直角只能位于 \(C'\)，即 \(B'\) 与 \(C'\) 等高：

\[
h_B=h_A. \tag{6}
\]

将式（2）、（4）代入：

\[
\frac{fl}{2(a-s)}=\frac{fl}{a}.
\]

约去非零的 \(fl\)：

\[
\frac1{2(a-s)}=\frac1a.
\]

交叉相乘：

\[
a=2(a-s),
\]

所以

\[
a=2s. \tag{7}
\]

由 \(a=u-f\)、\(s=\sqrt3l/2\) 得

\[
u-f=\sqrt3l,
\qquad
u=f+\sqrt3l. \tag{8}
\]

在条件（7）下，竖直直角边长度为

\[
h_A=\frac{fl}{2s}
=\frac{fl}{\sqrt3l}
=\frac f{\sqrt3}. \tag{9}
\]

水平直角边长度为

\[
\begin{aligned}
d
&=\frac{f^2s}{(2s)(s)}\\
&=\frac{f^2}{2s}\\
&=\frac{f^2}{\sqrt3l}.
\end{aligned}
\tag{10}
\]

因此两条直角边之比为

\[
\frac d{h_A}
=
\frac{f^2/(\sqrt3l)}{f/\sqrt3}
=\frac fl. \tag{11}
\]

要成为 \(30^\circ\)-\(60^\circ\)-\(90^\circ\) 三角形，\(d/h_A\) 只能等于
\(\sqrt3\) 或 \(1/\sqrt3\)。

\needspace{3\baselineskip}
\textbf{（三）第（1）（2）问：较大物距}

较大物距对应 \(l\) 较大。令竖直边为长直角边，即

\[
\frac d{h_A}=\frac1{\sqrt3}.
\]

由式（11）：

\[
\frac fl=\frac1{\sqrt3},
\]

所以

\[
\boxed{l_1=\sqrt3f}. \tag{12}
\]

代入式（8）：

\[
u_1=f+\sqrt3(\sqrt3f)=4f.
\]

故

\[
\boxed{u_1=4f}. \tag{13}
\]

复核像边长：

\[
h_A=\frac f{\sqrt3},
\qquad
d=\frac f3,
\qquad
\frac{h_A}{d}=\sqrt3.
\]

三角形确为 \(30^\circ\)-\(60^\circ\)-\(90^\circ\)。

\needspace{3\baselineskip}
\textbf{（四）第（3）（4）问：较小物距}

较小物距对应 \(l\) 较小。令水平边为长直角边，即

\[
\frac d{h_A}=\sqrt3.
\]

由式（11）：

\[
\frac fl=\sqrt3,
\]

所以

\[
\boxed{l_2=\frac f{\sqrt3}=\frac{\sqrt3}{3}f}. \tag{14}
\]

代入式（8）：

\[
u_2=f+\sqrt3\cdot\frac f{\sqrt3}=2f.
\]

故

\[
\boxed{u_2=2f}. \tag{15}
\]

此时

\[
h_A=\frac f{\sqrt3},
\qquad
d=f,
\qquad
\frac d{h_A}=\sqrt3,
\]

同样满足要求。

\needspace{3\baselineskip}
\textbf{（五）\(A'C'\) 为斜边时确定几何比例}

当像三角形的斜边垂直于主轴时，斜边只能是 \(A'C'\)，直角位于 \(B'\)。相对
\(A'\) 写出两个向量：

\[
\overrightarrow{B'A'}=(-d,h_B),
\qquad
\overrightarrow{B'C'}=(-d,h_B-h_A).
\]

直角条件为

\[
\overrightarrow{B'A'}\cdot\overrightarrow{B'C'}=0,
\]

即

\[
d^2+h_B(h_B-h_A)=0,
\]

所以

\[
d^2=h_B(h_A-h_B). \tag{16}
\]

对一个以 \(A'C'\) 为斜边的 \(30^\circ\)-\(60^\circ\)-\(90^\circ\)
三角形，直角顶点在斜边上的投影把斜边分成 \(1:3\) 的两段。因此

\[
\frac{h_B}{h_A}=\frac14
\quad\text{或}\quad
\frac34. \tag{17}
\]

而由式（2）、（4）直接得到

\[
\frac{h_B}{h_A}
=
\frac{a}{2(a-s)}. \tag{18}
\]

由于 \(a>s>0\)，式（18）恒大于 \(1/2\)，所以 \(1/4\)
分支不符合本题像点顺序，只能取

\[
\frac{h_B}{h_A}=\frac34. \tag{19}
\]

代入式（18）：

\[
\frac{a}{2(a-s)}=\frac34.
\]

交叉相乘：

\[
4a=6(a-s),
\]

整理得

\[
a=3s. \tag{20}
\]

\needspace{3\baselineskip}
\textbf{（六）第（5）（6）问：斜边垂直于主轴}

在 \(a=3s\) 时，\(a-s=2s\)。由式（2）、（5）：

\[
h_A=\frac{fl}{3s}
=\frac{2f}{3\sqrt3}, \tag{21}
\]

\[
d=\frac{f^2s}{(3s)(2s)}
=\frac{f^2}{6s}. \tag{22}
\]

对于 \(h_B/h_A=3/4\) 的分支，\(A'\) 处锐角为 \(30^\circ\)，所以

\[
d=\frac{\sqrt3}{4}h_A. \tag{23}
\]

将式（21）代入式（23）：

\[
d=\frac{\sqrt3}{4}\cdot\frac{2f}{3\sqrt3}
=\frac f6. \tag{24}
\]

比较式（22）与式（24）：

\[
\frac{f^2}{6s}=\frac f6,
\]

于是

\[
s=f. \tag{25}
\]

因为 \(s=\sqrt3l/2\)，

\[
\boxed{l_3=\frac{2f}{\sqrt3}=\frac{2\sqrt3}{3}f}. \tag{26}
\]

又由 \(a=3s=3f\)、\(u=a+f\)，

\[
\boxed{u_3=4f}. \tag{27}
\]

最后把结果代回完整坐标进行复核：

\[
h_A=\frac{2f}{3\sqrt3},
\qquad
h_B=\frac f{2\sqrt3},
\qquad
d=\frac f6.
\]

两条直角边长度分别为

\[
\begin{aligned}
|A'B'|
&=\sqrt{d^2+h_B^2}
=\frac f3,\\
|B'C'|
&=\sqrt{d^2+(h_A-h_B)^2}
=\frac f{3\sqrt3}.
\end{aligned}
\]

而

\[
|A'C'|=h_A=\frac{2f}{3\sqrt3}.
\]

因此

\[
|A'B'|:|B'C'|:|A'C'|
=\sqrt3:1:2,
\]

直角、边比与斜边方向三项条件全部满足。

\needspace{3\baselineskip}
\textbf{结果汇总与物理条件检查}

三组解为

\[
\boxed{
\begin{aligned}
u_1&=4f,& l_1&=\sqrt3f,\\
u_2&=2f,& l_2&=\frac f{\sqrt3},\\
u_3&=4f,& l_3&=\frac{2f}{\sqrt3}.
\end{aligned}}
\]

还需检查离透镜最近的点 \(B\) 是否仍在焦点外。其物距为 \(p_B=u-s\)：

\begin{itemize}
\tightlist
\item
  第一组：\(p_B=4f-\frac32f=\frac52f>f\)；
\item
  第二组：\(p_B=2f-\frac12f=\frac32f>f\)；
\item
  第三组：\(p_B=4f-f=3f>f\)。
\end{itemize}

所以三组均为实物成实像，所用成像公式和图中像点顺序自洽。

\needspace{3\baselineskip}
\textbf{逐点核对第三组解}

取 \(f=1\)，第三组物点条件为 \(u=4,l=2/\sqrt3\)。像点可取
\(A'=(0,0)\)、\(B'=(1/6,-1/(2\sqrt3))\)、\(C'=(0,-2/(3\sqrt3))\)。两条从
\(B'\) 出发的边的点积为 \(1/36-1/36=0\)；边长为
\(1/3,1/(3\sqrt3),2/(3\sqrt3)\)，比为
\(\sqrt3:1:2\)。图4-2保留三个物点的不同物距，用这些坐标检查直角，而不凭示意图上的角度读数。

\subsection{六、关键洞察}\label{ux516dux5173ux952eux6d1eux5bdf-3}

解析图的三组像三角形按逐点薄透镜坐标绘制，同一小图横纵比例一致。前两组的竖直像边都是直角边，分别较长和较短；第三组竖直像边为斜边，直角在B点。

第（5）问必须同时使用直角条件和 \(1:\sqrt3:2\) 边比。仅由
\(\overrightarrow{A'B'}\cdot\overrightarrow{C'B'}=0\)
只能得到一族直角三角形，不能唯一确定 \(u_3,l_3\)。

逐点坐标法还避免了``轴向放大率''被误用。公式 \(m_l=-m^2\)
只描述同一小邻域内沿主轴的微小长度变化；本题 \(A\) 与 \(B\)
的轴向间距可以与 \(f\)
同量级，且二者横向高度不同，直接拿一个局部轴向放大率处理整条边会丢失非线性。使用精确像距函数
\(q(p)=fp/(p-f)\) 再取差，既适用于有限间距，也能自动保留正确符号。

三组答案中 \(u_1=u_3=4f\) 并不表示像形相同：第一组边长
\(l_1=\sqrt3f\)，使 \(B'\) 与 \(C'\) 等高；第三组边长
\(l_3=2f/\sqrt3\)，使 \(B'\)
成为直角顶点。物距相同而物体尺度不同，改变了 \(B\) 的实际物距
\(u-s\)，因此像三角形的顶点角色随之改变。

还可用一个不含比例图误差的坐标表复查第（5）问。把所有长度除以 \(f\)，取

\[
A'=(0,0),\qquad
C'=\left(0,-\frac{2}{3\sqrt3}\right),\qquad
B'=\left(\frac16,-\frac1{2\sqrt3}\right).
\]

则

\[
\overrightarrow{B'A'}\cdot\overrightarrow{B'C'}
=\frac1{36}
+\frac1{2\sqrt3}\left(\frac1{2\sqrt3}-\frac{2}{3\sqrt3}\right)
=0.
\]

三边平方依次为
\(1/9\)、\(1/27\)、\(4/27\)，即长直角边、短直角边、斜边平方满足
\(3:1:4\)。这一代数验算同时排除了``图看起来近似直角''造成的误判，也说明
\(l_3=2f/\sqrt3\) 不是只由某一条角度关系偶然得到的。

\subsection{七、易错点}\label{ux4e03ux6613ux9519ux70b9-3}

\begin{enumerate}
\def\labelenumi{\arabic{enumi}.}
\tightlist
\item
  把原图中垂直主轴的边误认为 \(BC\)，从一开始就建立错误坐标。
\item
  对三个顶点使用同一个放大率。\(B\) 的物距是 \(u-\sqrt3l/2\)，与 \(A,C\)
  不同。
\item
  忽略倒立像，只写像高绝对值后又在向量点积中丢失上下关系。
\item
  第（1）（3）问只得到 \(u=f+\sqrt3l\)，没有继续使用直角边比求 \(l\)。
\item
  第（5）问只列垂直条件，没有利用 \(30^\circ\)-\(60^\circ\)-\(90^\circ\)
  的投影比例。
\item
  得到代数根后未检查 \(u-s>f\)，从而可能把虚像分支混入实像图。
\end{enumerate}

验算时要求每组答案都过三关：最近物点 \(B\) 的物距大于 \(f\)；用
\(h_A,h_B,d\) 算出的两边点积为零；三条边长度比确为
\(1:\sqrt3:2\)。特别是第（5）问，若只得到 \(u_3=4f\) 而未确定
\(l_3\)，或给出的 \(l_3\)
不能同时通过点积和边比检验，均属于条件使用不全。

\subsection{八、课堂讲解建议}\label{ux516bux8bfeux5802ux8bb2ux89e3ux5efaux8bae-3}

先让学生在图上标出 \(A,C\) 物距相同、\(B\) 更靠近透镜，再只计算
\(q_A,q_B\) 的大小关系。随后把像方抽象成一个竖直线段 \(A'C'\) 和线外一点
\(B'\)，先完成纯几何分类，最后代入光学公式。这样可以把``看图困难''和``代数困难''分开处理。

讲第（5）问时，可先在任意 \(30^\circ\)-\(60^\circ\)-\(90^\circ\)
三角形中向斜边作高，让学生自行证明斜边被分成 \(1:3\)，再把这一比例对应到
\(h_B/h_A\)。

板书第一行只画原图并写 \(p_A=p_C=u\)、\(p_B=u-s\)，要求学生先判断
\(q_B>q_A\)。第二行再写 \(h_A,h_B,d\)
三个量，后续六问全部只在这三个量上运算。这样既能控制代数长度，也能不断提醒学生每个式子对应图中的哪一段。

三个结果求出后，不应立即结束。让学生逐组计算
\(p_B=u-s\)，检查最近物点仍在焦点外；再把 \(h_A,h_B,d\)
代回，直接算三条像边长度。第（5）问尤其要用点积和 \(1:\sqrt3:2\)
两次检验。课堂上若只核对最终数值，学生很难发现``直角成立但角度不是
\(30^\circ,60^\circ\)''这一类隐蔽错误。

\subsection{九、变式训练}\label{ux4e5dux53d8ux5f0fux8badux7ec3-3}

\begin{enumerate}
\def\labelenumi{\arabic{enumi}.}
\tightlist
\item
  将正三角形换成顶角为 \(2\alpha\) 的等腰三角形，建立 \(h_A,h_B,d\)
  的一般表达式。
\item
  若要求像三角形为等腰直角三角形，求 \(u\) 与 \(l\) 的关系。
\item
  保持 \(l\) 固定，讨论 \(u\)
  从焦点外向远处移动时，像三角形三个内角如何变化。
\item
  允许 \(B\) 位于焦点内，重新讨论虚像与实像混合时的几何分支。
\end{enumerate}

\textbf{变式解析要点：} 对顶角为 \(2\alpha\) 的等腰三角形，只需把

\[
s=d\cos\alpha,\qquad y_B=d\sin\alpha,\qquad AC=2d\sin\alpha
\]

这里 \(d\) 为两条等腰边长、\(2\alpha\) 为顶角，不能把原题边长 \(l\)
直接沿用。代入逐点成像公式，\(A,C\) 的像仍具有相同像距。等腰直角像可在
\(A'C'\) 为直角边时令 \(d=h_A\)，结合 \(h_B=h_A\) 求解；若 \(A'C'\)
为斜边，则令
\(h_B/h_A=1/2\)，但还要检查式（18）能否达到该值。研究虚像分支时最好改用带符号像距
\(q\)，并始终用坐标而非``倒立像''口语判断上下；最后通过光心射线是否连续、各边点积是否为零来验算。

\newpage

\section{第 5 题}\label{ux7b2c-5-ux9898}

如图所示，有一惯性系 \(S\)，另一惯性系 \(S'\) 相对 \(S\) 以
\(\beta c\,\bm i\) 运动，\(S''\) 相对 \(S\) 以
\(\beta c(\cos\theta\,\bm i+\sin\theta\,\bm j)\) 运动，设
\(0\le\beta<1\)；比较相邻随动系时取
\(|\theta|\ll1\)。这些惯性系采用标准对时方式，即设三个惯性系的原点同时相遇，并各自在自己的参考系中将各处时钟同时调为
\(0\)。

（1）将 \(x'',y'',t''\) 用 \(x',y',t'\) 表示。为简化，可引入
\(\gamma=\dfrac1{\sqrt{1-\beta^2}}\)。（提示：可作坐标变换
\(\begin{cases}\xi=\cos\theta\,x+\sin\theta\,y\\ \eta=-\sin\theta\,x+\cos\theta\,y\end{cases}\)，\(\begin{cases}\xi''=\cos\theta\,x''+\sin\theta\,y''\\ \eta''=-\sin\theta\,x''+\cos\theta\,y''\end{cases}\)，则
\((\xi,\eta,t)\) 与 \((\xi'',\eta'',t'')\) 满足常规形式的洛伦兹变换。）

（2）在 \(t'=t''=0\) 的同时面上，\(S''\) 系的坐标架相对 \(S'\)
系转过一小角度 \(d\varphi\)（逆时针为正）。设 \(S',S''\)
分别为一粒子在时刻 \(t\) 和 \(t+dt\)
的``随动系''，即粒子的``随动系''坐标架以 \(\bm\omega_t\) 旋转，且
\(\bm\omega_t\) 的正方向与 \(d\varphi\) 一致。用粒子的 \(\bm a,\bm v\)
表示 \(\bm\omega_t\)，已知 \(|\bm v|\ll c\) 且
\(\bm a\perp\bm v\)。这一现象称为托马斯旋进，对电子的自旋---轨道相互作用有重要影响。

\begin{figure}[H]
\centering
\includegraphics[width=0.7\linewidth,height=\textheight,keepaspectratio,alt={图5-1　惯性系 S\textquotesingle{} 与 S\textquotesingle\textquotesingle{} 相对同一惯性系 S 的等速、不同方向运动。}]{../figures/question-05.pdf}
\caption*{图5-1　惯性系 \(S'\) 与 \(S''\) 相对同一惯性系 \(S\)
的等速、不同方向运动。}
\end{figure}

\subsection{一、题型定位}\label{ux4e00ux9898ux578bux5b9aux4f4d-4}

两次速度都相对实验室系 \(S\)
给出。先复合标准洛伦兹变换，再取速度方向的小角变化，才能读出相邻随动坐标轴的转角及其方向约定。

\subsection{二、知识点}\label{ux4e8cux77e5ux8bc6ux70b9-4}

\begin{enumerate}
\def\labelenumi{\arabic{enumi}.}
\tightlist
\item
  沿任意单位方向 \(\bm n\) 的标准洛伦兹变换。
\item
  平面坐标的``旋转---沿轴 boost---旋回''分解。
\item
  由 \(S\to S'\) 的逆变换把 \(x,y,t\) 消去。
\item
  同时性条件 \(t'=t''=0\) 与空间轴转角的读取。
\item
  低速展开 \(\gamma-1\simeq v^2/(2c^2)\)。
\item
  \(\bm a\perp\bm v\) 时速度方向转率 \(d\theta/dt=a/v\)，以及向量方向
  \(\bm v\times\bm a\)。
\end{enumerate}

\subsection{三、物理图像}\label{ux4e09ux7269ux7406ux56feux50cf-4}

在时刻 \(t\)，粒子速度沿 \(x\) 轴，其标准随动系是 \(S'\)；经过 \(dt\)
后，速度方向在实验室系中转过 \(d\theta\)，新的标准随动系是
\(S''\)。两套坐标都由实验室系按``无额外空间转动''的标准 boost
定义，但把二者直接比较时，空间轴之间出现

\[
d\varphi=(\gamma-1)d\theta.
\]

这个小转角是非共线 boost 不可交换的结果。题目规定的正方向与 \(d\varphi\)
相同，因此低速时 \(\bm\omega_t\) 沿 \(\bm v\times\bm a\)。

\subsection{四、解题入口}\label{ux56dbux89e3ux9898ux5165ux53e3-4}

先写 \(S\to S'\) 的标准变换及其逆变换。再把 \(S''\) 相对 \(S\)
的速度方向记为 \(\bm n=(\cos\theta,\sin\theta)\)，沿 \(\bm n\) 与垂直
\(\bm n\) 方向分解坐标，得到 \(x'',y'',t''\) 关于 \(x,y,t\)
的式子。最后代入逆变换
\(x=\gamma(x'+\beta ct')\)、\(ct=\gamma(ct'+\beta x')\)、\(y=y'\)。

第（2）问只保留 \(\theta\) 的一阶小量，在 \(t'=t''=0\) 时比较

\[
x''=x'+d\varphi\,y',\qquad
y''=y'-d\varphi\,x',
\]

即可读出 \(d\varphi\)；再用 \(d\theta/dt=a/v\) 与
\(\gamma-1\simeq v^2/(2c^2)\) 转成题目所要的向量角速度。

\subsection{五、解析}\label{ux4e94ux89e3ux6790-4}

\begin{figure}[H]
\centering
\includegraphics[width=0.9\linewidth,height=\textheight,keepaspectratio,alt={图5-2　相邻时刻的速度方向变化、随动系轴转角与加速度方向。}]{../figures/solution-05.pdf}
\caption*{图5-2　相邻时刻的速度方向变化、随动系轴转角与加速度方向。}
\end{figure}

图5-2用于逐项读取相对变换矩阵，两次标准推动的速度都相对于S给定；空间子块不直接等同旋转矩阵，后面另检查洛伦兹度规。

图 5-2
以实验室系观察粒子在相邻时刻的速度：\(\bm v(t+dt)-\bm v(t)=\bm a\,dt\)，速度方向转过
\(d\theta\)；由两套标准随动坐标的空间部分读取坐标架转角 \(d\varphi\)。

\needspace{3\baselineskip}
\textbf{（一）第（1）问：分别从 \(S\) 建立两个标准洛伦兹变换}

\(S'\) 相对 \(S\) 沿 \(x\) 轴以 \(\beta c\) 运动，标准洛伦兹变换为

\[
x'=\gamma(x-\beta ct),\qquad y'=y,\qquad
ct'=\gamma(ct-\beta x). \tag{1}
\]

\(S''\) 相对 \(S\) 的速度方向为

\[
\bm n=(\cos\theta,\sin\theta).
\]

先在 \(S\) 系中把坐标分解到平行、垂直于 \(\bm n\) 的方向：

\[
\xi=\cos\theta\,x+\sin\theta\,y,\qquad
\eta=-\sin\theta\,x+\cos\theta\,y. \tag{2}
\]

于是 \(S''\) 相对 \(S\) 沿 \(\xi\) 轴以 \(\beta c\) 运动。对
\((\xi,\eta,t)\) 应用标准洛伦兹变换：

\[
\xi''=\gamma(\xi-\beta ct),\qquad
\eta''=\eta,\qquad
ct''=\gamma(ct-\beta\xi). \tag{3}
\]

再把 \(S''\) 中沿速度方向与垂直速度方向的坐标旋回 \(x'',y''\)：

\[
x''=\cos\theta\,\xi''-\sin\theta\,\eta'',\qquad
y''=\sin\theta\,\xi''+\cos\theta\,\eta''. \tag{4}
\]

为缩短书写，记

\[
C=\cos\theta,\qquad S_\theta=\sin\theta.
\]

把式（2）（3）代入式（4），先得到 \(S''\) 关于实验室系 \(S\) 的变换：

\[
\begin{aligned}
ct''&=\gamma\left(ct-\beta Cx-\beta S_\theta y\right),\\
x''&=(\gamma C^2+S_\theta^2)x
+(\gamma-1)CS_\theta y-\gamma\beta C\,ct,\\
y''&=(\gamma-1)CS_\theta x
+(\gamma S_\theta^2+C^2)y-\gamma\beta S_\theta\,ct.
\end{aligned}
\tag{5}
\]

现在由式（1）的逆变换消去 \(S\) 的坐标：

\[
x=\gamma(x'+\beta ct'),\qquad
ct=\gamma(ct'+\beta x'),\qquad
y=y'. \tag{6}
\]

代入式（5），得到题目要求的完整关系：

\[
\begin{aligned}
ct''
&=\gamma^2(1-\beta^2C)\,ct'
+\gamma^2\beta(1-C)\,x'
-\gamma\beta S_\theta\,y',\\[3pt]
x''
&=(\gamma C^2+S_\theta^2)\gamma(x'+\beta ct')
+(\gamma-1)CS_\theta\,y'\\
&\quad-\gamma^2\beta C(ct'+\beta x'),\\[3pt]
y''
&=(\gamma-1)CS_\theta\,\gamma(x'+\beta ct')
+(\gamma S_\theta^2+C^2)y'\\
&\quad-\gamma^2\beta S_\theta(ct'+\beta x').
\end{aligned}
\tag{7}
\]

这里保留因式形式，是为了让每一项都能追溯到式（5）的空间旋转或
boost。可用两个极限复核：

\begin{itemize}
\tightlist
\item
  \(\theta=0\) 时 \(S''\) 与 \(S'\) 完全相同，式（7）给出
  \(x''=x',y''=y',t''=t'\)；
\item
  \(\beta=0\) 时 \(\gamma=1\)，三个系都静止，同样退化为恒等变换。
\end{itemize}

\needspace{3\baselineskip}
\textbf{（二）第（2）问：从空间变换中读出小转角}

对相邻两个随动系，速度方向只改变一个小角 \(d\theta\)。在式（7）中取

\[
C=1+O(d\theta^2),\qquad
S_\theta=d\theta+O(d\theta^3).
\]

只保留 \(d\theta\) 的一阶项，逐项化简：

\[
\begin{aligned}
ct''&=ct'-\gamma\beta\,d\theta\,y',\\
x''&=x'+(\gamma-1)d\theta\,y',\\
y''&=y'-(\gamma-1)d\theta\,x'
-\gamma\beta\,d\theta\,ct'.
\end{aligned}
\tag{8}
\]

式（8）的纯空间部分与小角坐标旋转

\[
\begin{pmatrix}x''\\y''\end{pmatrix}
=
\begin{pmatrix}1&d\varphi\\-d\varphi&1\end{pmatrix}
\begin{pmatrix}x'\\y'\end{pmatrix}
\]

同形，因此在题目规定的正方向下

\[
\boxed{d\varphi=(\gamma-1)d\theta}. \tag{9}
\]

时间式中的 \(-\gamma\beta\,d\theta\,y'\)
表明两个随动系的同时面并不重合；这正是不能把问题当作普通欧氏旋转的原因。读取空间轴方向时必须同时保留式（8）的时间关系，不能先假定所有空间点都满足
\(t'=t''=0\)。

\needspace{3\baselineskip}
\textbf{（三）用 \(\bm a,\bm v\) 表示角速度}

因为 \(\bm a\perp\bm v\)，在 \(dt\)
内速度大小不变，只有方向改变。速度矢量端点走过的小弧长为
\(a\,dt\)，而速度圆的半径为 \(v\)，故

\[
\frac{d\theta}{dt}=\frac av. \tag{10}
\]

由式（9）（10），随动坐标架按题目定义的旋进角速度大小为

\[
\omega_t=(\gamma-1)\frac av. \tag{11}
\]

方向由速度方向转向加速度方向的右手定则确定，即沿
\(\bm v\times\bm a\)。于是向量形式为

\[
\bm\omega_t
=\frac{\gamma-1}{v^2}\,\bm v\times\bm a. \tag{12}
\]

题设 \(v\ll c\)，所以

\[
\gamma
=\frac1{\sqrt{1-v^2/c^2}}
=1+\frac12\frac{v^2}{c^2}
+O\!\left(\frac{v^4}{c^4}\right).
\]

代入式（12）：

\[
\boxed{\bm\omega_t
\simeq\frac{1}{2c^2}\,\bm v\times\bm a}. \tag{13}
\]

量纲检查：\(va/c^2\) 的单位为 \(\mathrm{s^{-1}}\)，正是角速度。且
\(v/c\to0\) 时结果按 \(v^2/c^2\) 的阶数消失，不会出现一阶相对论效应。

\needspace{3\baselineskip}
\textbf{（四）变换结构与符号复核}

二维洛伦兹变换最稳妥的写法是先把空间位置分解为平行、垂直于速度方向的分量：

\[
\bm r'=\bm r_\perp+\gamma(\bm r_\parallel-\bm vt),
\qquad
ct'=\gamma\left(ct-\frac{\bm v\cdot\bm r}{c}\right).
\]

取 \(\bm n=(\cos\theta,\sin\theta)\)，就有
\(\bm r_\parallel=(\bm n\cdot\bm r)\bm n\)。展开后，\(x''\) 与 \(y''\)
的交叉项必同时含
\((\gamma-1)\sin\theta\cos\theta\)。这个结构检验比逐项背公式可靠。

还要区分两种常见符号约定。本题把``\(S''\) 坐标架相对 \(S'\)
的转角''定义为正，所以式（13）沿
\(\bm v\times\bm a\)。若改问``无力矩自旋轴相对轨道方向的进动''，常用的托马斯角速度写成沿
\(\bm a\times\bm v\)，即整体差一个负号。批改时必须先核对被比较的坐标架和正方向。

\needspace{3\baselineskip}
\textbf{时空间隔与转角约定}

令 \(X=(ct,x,y)\)，标准变换矩阵为 \(\Lambda(\bm\beta)\)，则两系间矩阵是
\(\Lambda(\bm\beta'')\Lambda(-\bm\beta')\)。每个因子满足
\(\Lambda^{\mathsf T}\eta\Lambda=\eta\)，其中
\(\eta=\operatorname{diag}(1,-1,-1)\)，所以式（7）也必须保持
\((ct)^2-x^2-y^2\)。这给出独立于展开步骤的检查。

本文正转角表示坐标轴沿 \(\bm v\times\bm a\) 的方向转动，对应被动坐标式
\(x''=x'+d\varphi y',y''=y'-d\varphi x'\)。若讨论保持无转矩的自旋相对这些坐标轴的转动，符号约定改变；不能把两种角速度混写。

\subsection{六、关键洞察}\label{ux516dux5173ux952eux6d1eux5bdf-4}

解析图用于逐项读取相对变换矩阵，两次标准推动的速度都相对于S给定；空间子块不直接等同旋转矩阵，后面另检查洛伦兹度规。

第二，托马斯转角来自``同时性与空间方向不可分割''。式（8）不仅有空间交叉项，还有

\[
ct''=ct'-\gamma\beta\,d\theta\,y'.
\]

因此不能把 \(t'=t''=0\)
当作覆盖整个平面的共同同时面。真正可读出的轴转角来自完整四维变换的一阶空间部分，其系数是
\(\gamma-1\)。

第三，题目定义的是相邻两套标准随动坐标架之间的正向转角，所以

\[
d\varphi=(\gamma-1)d\theta,
\qquad
\bm\omega_t\simeq\frac{\bm v\times\bm a}{2c^2}.
\]

教科书中常见的``自旋相对轨道的托马斯进动''使用相反比较顺序，写成
\(\bm a\times\bm v/(2c^2)\)。这不是物理矛盾，而是被比较坐标架和正方向不同；课堂上必须把这一点说清楚。

\subsection{七、易错点}\label{ux4e03ux6613ux9519ux70b9-4}

\begin{enumerate}
\def\labelenumi{\arabic{enumi}.}
\tightlist
\item
  把``\(S''\) 相对 \(S\)''误抄为``\(S''\) 相对 \(S'\)''。
\item
  第（1）问只给出 \(x'',y''\) 关于 \(x,y,t\)
  的式子，忘记继续用逆变换消去实验室坐标。
\item
  旋转坐标时只``转过去''而没有``旋回来''，导致 \(x'',y''\)
  的交叉项错误。
\item
  把 \(\gamma^2\beta^2=\gamma^2-1\) 漏用或误写，导致 \(\theta=0\)
  时不能退化为恒等变换。
\item
  小角展开时保留 \(\sin\theta\simeq\theta\)，却又错误保留
  \(1-\cos\theta\) 的一阶项；后者实际是二阶小量。
\item
  直接令所有空间点同时满足 \(t'=t''=0\)，忽略两个同时面的差别。
\item
  把 \(d\varphi\) 与实验室中的速度方向转角 \(d\theta\) 混为一谈。
\item
  低速展开漏掉因子 \(1/2\)，或把 \(va/c^2\) 写成 \(a/(vc^2)\)。
\item
  未按题设正方向判断叉积次序。
\end{enumerate}

批改时先查题意：两套速度是否都相对 \(S\)。再查三个极限：\(\theta=0\)
应为恒等变换，\(\beta=0\) 应为恒等变换，\(v/c\to0\) 时旋进应按
\(v^2/c^2\) 消失。三个检查中任一个不通过，都说明变换或符号存在根本问题。

\subsection{八、课堂讲解建议}\label{ux516bux8bfeux5802ux8bb2ux89e3ux5efaux8bae-4}

先在同一张速度图上画 \(\bm v(t)\) 与
\(\bm v(t+dt)\)，明确二者都由实验室系测量。随后分两列写变换：左列是沿
\(x\) 的 \(S\to S'\)，右列是``旋转---boost---旋回''的
\(S\to S''\)。直到两列都完成后，才用式（6）消去
\(S\)；这样可以有效避免把两个 boost 错当成顺次作用。

处理长式（7）时不必一开始完全展开。让学生保留
\((\gamma C^2+S_\theta^2)\gamma(x'+\beta ct')\) 这样的因式，并立即做
\(\theta=0\) 检查。通过后再取 \(d\theta\)
一阶量，学生会清楚看到空间交叉项恰为 \(\pm(\gamma-1)d\theta\)。

最后用速度矢量端点的小圆弧解释
\(d\theta/dt=a/v\)。先写大小，再用右手定则确定
\(\bm v\times\bm a\)；然后另列一行说明常见自旋进动约定为何会改成
\(\bm a\times\bm v\)。把``大小计算''和``符号约定''分开，可以减少学生因一个负号而否定整段推导。

\subsection{九、变式训练}\label{ux4e5dux53d8ux5f0fux8badux7ec3-4}

\begin{enumerate}
\def\labelenumi{\arabic{enumi}.}
\tightlist
\item
  不作小角近似，使用式（7）验证 \(\theta=\pi/2\) 时的坐标变换。
\item
  粒子以速度 \(v\) 作半径为 \(R\)
  的匀速圆周运动，求一周内随动坐标架累积的低速托马斯转角。
\item
  若 \(\bm a\parallel\bm v\)，判断是否出现托马斯转角。
\item
  把低速近似去掉，保留 \(\gamma\)，写出 \(\bm\omega_t\)
  的精确形式并与常用自旋进动公式比较。
\end{enumerate}

\textbf{变式解析要点：} 第 1 项取 \(C=0,S_\theta=1\)
直接代入式（7），并检查 \(S''\) 的 boost 确实沿实验室 \(y\) 轴。第 2
项中 \(a=v^2/R\)、轨道角速度为 \(v/R\)，故

\[
\omega_t\simeq\frac{v^3}{2Rc^2},
\qquad
\Delta\varphi=\omega_t\frac{2\pi R}{v}
=\pi\frac{v^2}{c^2}.
\]

方向按本题定义沿轨道转动方向；若改讨论自旋相对轨道的进动，方向相反。第 3
项中速度方向不变，\(d\theta=0\)，所以没有由非共线 boost 产生的转角。第 4
项按本题轴转角定义有

\[
\bm\omega_t=\frac{\gamma-1}{v^2}\bm v\times\bm a
\]

（这里假设 \(\bm a\perp\bm v\)）；常用自旋进动公式则因比较顺序相反写成

\[
\bm\Omega_T=\frac{\gamma^2}{\gamma+1}\frac{\bm a\times\bm v}{c^2},
\]

并可用 \(\gamma^2\beta^2=(\gamma-1)(\gamma+1)\) 验证两者大小相同。

\newpage

\section{第 6 题}\label{ux7b2c-6-ux9898}

（1）\(n\) 摩尔理想气体的摩尔定容热容为 \(C_V\)。已知在
\(p=p_0\)、\(T=T_0\) 时熵为 \(S_0\)，求 \(S(p,T)\)。

（2）化学势定义为

\[
\mu=\frac{U+pV-TS}{n}.
\]

取 \(U|_{T=0}=0\)，求 \(\mu(p,T)\)。

（3）对反应
\(n\to p+e^-+\bar\nu\)，平衡条件为两侧反应物化学势之和相等。取
\(\mu_e\simeq N_A m_ec^2\)、\(\mu_{\bar\nu}\simeq0\)；质子、中子均视作理想气体，二者的
\(C_V,p_0,T_0,S_0/n\) 相同。求平衡时质子分压比 \(x_p=p_p/p\)。

\begin{figure}[H]
\centering
\includegraphics[width=0.66\linewidth,height=\textheight,keepaspectratio,alt={图6-1　理想气体状态与反应组分}]{../figures/question-06.pdf}
\caption*{图6-1　理想气体状态与反应组分}
\end{figure}

\subsection{一、题型定位}\label{ux4e00ux9898ux578bux5b9aux4f4d-5}

先由理想气体的全微分求熵，再除以物质的量求摩尔化学势。反应平衡式必须统一每摩尔与每粒子的能量，指数才能无量纲。

\subsection{二、知识点}\label{ux4e8cux77e5ux8bc6ux70b9-5}

\begin{enumerate}
\def\labelenumi{\arabic{enumi}.}
\tightlist
\item
  可逆过程有 \(T\,dS=dU+p\,dV\)。
\item
  理想气体满足 \(dU=nC_VdT\)、\(pV=nRT\)。
\item
  \(C_p=C_V+R\)。
\item
  摩尔吉布斯自由能等于化学势：\(\mu=(U+pV-TS)/n\)。
\item
  化学平衡条件为化学计量数加权后的化学势之和为零。
\item
  同温理想混合气体中，组分化学势含 \(RT\ln p_i\)。
\end{enumerate}

\subsection{三、物理图像}\label{ux4e09ux7269ux7406ux56feux50cf-5}

从参考态 \((p_0,T_0)\) 到任意态
\((p,T)\)，可选一条``先等压改变温度，再等温改变压强''的可逆路径。熵是状态函数，因此积分值与路径无关。第三问中，中子与质子的热学参数相同，二者化学势的公共部分相消，只留下
\(RT\ln(p_n/p_p)\)；电子静能给出偏向中子一侧的能量代价。

\subsection{四、解题入口}\label{ux56dbux89e3ux9898ux5165ux53e3-5}

把 \(V=nRT/p\) 微分，代入

\[
dS=\frac{nC_VdT+p\,dV}{T},
\]

化为只含 \(dT/T\) 和 \(dp/p\) 的全微分。得到 \(S\)
后直接代入化学势定义。第三问写

\[
\mu_n=\mu_p+\mu_e+\mu_{\bar\nu},
\]

并使用 \(p_n=(1-x_p)p\)、\(p_p=x_pp\)。

\subsection{五、解析}\label{ux4e94ux89e3ux6790-5}

\begin{figure}[H]
\centering
\includegraphics[width=0.9\linewidth,height=\textheight,keepaspectratio,alt={图6-2　参考态积分路径与化学势平衡}]{../figures/solution-06.pdf}
\caption*{图6-2　参考态积分路径与化学势平衡}
\end{figure}

图6-2由总量自由能转到每摩尔化学势，反应静能也乘阿伏伽德罗常数。这样
指数为每摩尔能量与 \(RT\) 的比值。等价的两种写法是

\[\frac{N_A m_ec^2}{RT}=\frac{m_ec^2}{k_BT}.\]

图 6-2 左图给出熵积分路径，右图表示中子衰变反应两侧的化学势守恒。

\needspace{3\baselineskip}
\textbf{（一）熵的全微分}

由热力学第一、第二定律的可逆形式

\[
T\,dS=dU+p\,dV. \tag{1}
\]

理想气体内能只依赖温度：

\[
dU=nC_VdT. \tag{2}
\]

由状态方程 \(V=nRT/p\) 得

\[
\begin{aligned}
dV
&=nR\,d\!\left(\frac{T}{p}\right)\\
&=nR\left(\frac{dT}{p}-\frac{T}{p^2}dp\right).
\end{aligned} \tag{3}
\]

将式（2）、（3）代入式（1）：

\[
\begin{aligned}
T\,dS
&=nC_VdT+p\,nR\left(\frac{dT}{p}-\frac{T}{p^2}dp\right)\\
&=n(C_V+R)dT-nRT\frac{dp}{p}.
\end{aligned}
\]

两边除以 \(T\)：

\[
dS=n(C_V+R)\frac{dT}{T}-nR\frac{dp}{p}. \tag{4}
\]

从 \((p_0,T_0,S_0)\) 积分至 \((p,T,S)\)：

\[
S-S_0=n(C_V+R)\ln\frac{T}{T_0}-nR\ln\frac{p}{p_0}.
\]

所以

\[
\boxed{S(p,T)=S_0+n(C_V+R)\ln\frac{T}{T_0}-nR\ln\frac{p}{p_0}}. \tag{5}
\]

\needspace{3\baselineskip}
\textbf{（二）化学势}

由 \(U(0)=0\) 与 \(dU=nC_VdT\)：

\[
U=nC_VT. \tag{6}
\]

又有

\[
pV=nRT. \tag{7}
\]

代入定义：

\[
\begin{aligned}
\mu
&=\frac{nC_VT+nRT-TS}{n}\\
&=(C_V+R)T-\frac{T}{n}S.
\end{aligned} \tag{8}
\]

再代入式（5）：

\[
\begin{aligned}
\mu(p,T)
={}&\left(C_V+R-\frac{S_0}{n}\right)T\\
&-(C_V+R)T\ln\frac{T}{T_0}
+RT\ln\frac{p}{p_0}.
\end{aligned}
\]

故

\[
\boxed{\mu(p,T)=\left(C_V+R-\frac{S_0}{n}\right)T-(C_V+R)T\ln\frac{T}{T_0}+RT\ln\frac{p}{p_0}}. \tag{9}
\]

量纲检查：每一项都是``每摩尔能量''，单位为 \(\mathrm{J\,mol^{-1}}\)。

\needspace{3\baselineskip}
\textbf{（三）反应平衡与质子分压比}

平衡条件为

\[
\mu_n=\mu_p+\mu_e+\mu_{\bar\nu}. \tag{10}
\]

按题设

\[
\mu_e\simeq N_A m_ec^2,\qquad \mu_{\bar\nu}\simeq0. \tag{11}
\]

中子与质子的温度及参考态参数相同，因此式（9）中除压强对数项外的部分相消：

\[
\mu_n-\mu_p=RT\ln\frac{p_n}{p_p}. \tag{12}
\]

由式（10）---（12）：

\[
RT\ln\frac{p_n}{p_p}=N_A m_ec^2. \tag{13}
\]

指数化：

\[
\frac{p_n}{p_p}=\exp\!\left(\frac{N_A m_ec^2}{RT}\right). \tag{14}
\]

定义 \(x_p=p_p/p\)，则

\[
p_p=x_pp,\qquad p_n=(1-x_p)p. \tag{15}
\]

代入式（14）：

\[
\frac{1-x_p}{x_p}=\exp\!\left(\frac{N_A m_ec^2}{RT}\right).
\]

移项求解：

\[
\boxed{x_p=\frac{1}{1+\exp\!\left(N_A m_ec^2/RT\right)}}. \tag{16}
\]

当 \(T\to0\) 时指数趋于无穷，\(x_p\to0\)；当 \(T\to\infty\) 时指数趋于
\(1\)，\(x_p\to1/2\)，这是本题简化化学势模型的形式极限；常热容经典气体的熵式不能一直外推到零温，真实反应平衡还需要各物种的统计权重和标准化学势。

\needspace{3\baselineskip}
\textbf{热力学微分与化学平衡的复核}

第一问可从基本关系逐项核对。对理想气体

\[
dU=nC_V\,dT,
\qquad
p\,dV=nR\,dT-nRT\frac{dp}{p}.
\]

代入 \(TdS=dU+p\,dV\) 得

\[
dS=n(C_V+R)\frac{dT}{T}-nR\frac{dp}{p}.
\]

右侧是全微分，分别从 \((p_0,T_0)\) 积到
\((p,T)\)，便得到式（5）。对结果求偏导可反向检验：

\[
\left(\frac{\partial S}{\partial T}\right)_p
=\frac{n(C_V+R)}{T},
\qquad
\left(\frac{\partial S}{\partial p}\right)_T
=-\frac{nR}{p}.
\]

前者正是定压热容
\(nC_p/T\)，后者与理想气体热膨胀关系一致。对数的自变量均为比值
\(T/T_0\)、\(p/p_0\)，所以无量纲；这是检查熵表达式时必须满足的条件。

化学势由摩尔 Gibbs 自由能给出：

\[
\mu=\frac{G}{n}=\frac{U+pV-TS}{n}.
\]

把 \(U=nC_VT\)、\(pV=nRT\) 与 \(S(p,T)\) 全部代入时，要把参考熵项写成
\(S_0/n\)，因为 \(S_0\) 是总熵而 \(\mu\) 是每摩尔量。代入后对 \(p\)
求偏导应满足

\[
\left(\frac{\partial\mu}{\partial p}\right)_T
=\frac{RT}{p}=V_m,
\]

这为式（9）提供独立校验。

第三问能消去共同的温度项和参考态项，是因为中子、质子采用相同的理想气体参数；若二者参考化学势不同，方程右侧还应保留常量差。现有假设下

\[
RT\ln\frac{p_n}{p_p}=N_A m_ec^2
\]

这里 \(m_e\) 是单个电子质量、\(N_A\) 是阿伏伽德罗常数，故 \(N_A m_ec^2\)
为每摩尔能量，分子、分母均采用每摩尔量，所得指数无量纲。原扫描和原整理稿的
\(m_ec^2/(RT)\) 混用了单位，不能只在注释中说明而保留错误指数。

\needspace{3\baselineskip}
\textbf{模型和低温极限的边界}

题设人为令中子、质子的参考化学势相同，本文比例是这一简化理想气体模型的结果，并未包含真实两种粒子的静能差与统计简并。\(U=nC_VT\)
是给定常热容模型下的形式能量基准；熵中的 \(\ln T\)
表明该经典模型不能一直外推到绝对零度。\(x_p\to0\) 和 \(1/2\)
只检验给定数学公式的极限。

\subsection{六、关键洞察}\label{ux516dux5173ux952eux6d1eux5bdf-5}

解析图由总量自由能转到每摩尔化学势，反应静能也乘阿伏伽德罗常数。这样
指数为每摩尔能量与 \(RT\) 的比值。等价的两种写法是

\[\frac{N_A m_ec^2}{RT}=\frac{m_ec^2}{k_BT}.\]

\subsection{七、易错点}\label{ux4e03ux6613ux9519ux70b9-5}

\begin{enumerate}
\def\labelenumi{\arabic{enumi}.}
\tightlist
\item
  把 \(C_V\) 与 \(C_p=C_V+R\) 混淆，漏掉 \(p\,dV\) 中的 \(nR\,dT\)。
\item
  对 \(V=nRT/p\) 微分时漏掉 \(-nRTdp/p^2\)。
\item
  将总熵 \(S_0\) 误当摩尔熵，导致化学势中少除一个 \(n\)。
\item
  在化学平衡式中把电子静能写到错误一侧。
\item
  误写 \(p_n=x_pp\)；正确的是 \(p_n=(1-x_p)p\)。
\end{enumerate}

验算清单包括：对数内部必须无量纲；\(dS\) 对 \(T,p\)
的偏导应还原基本微分；\(\partial\mu/\partial p=RT/p\)；化学平衡指数在
\(T\to0\) 与 \(T\to\infty\)
时给出合理极限。还要统一``每摩尔''和``每粒子''的能量单位，每摩尔静能是
\(N_A m_ec^2\)；若按单粒子计，则用 \(m_ec^2\) 和 \(k_B\)。

\subsection{八、课堂讲解建议}\label{ux516bux8bfeux5802ux8bb2ux89e3ux5efaux8bae-5}

先要求学生独立完成式（3）的全微分，再用等容---等温路径分别积分复核式（5）。第三问可先隐藏式（16），让学生用低温极限判断指数应位于分母还是分子。

课堂上可把 \(TdS=dU+p\,dV\) 的每一项都标注单位，并在代入 \(V=nRT/p\)
后先不约分，让学生看清 \(nC_VdT+nRdT\) 为什么合成
\(nC_pdT\)。积分完成后，通过对 \(T,p\)
求偏导反查微分式，借此训练``从结果回到微分''的状态函数校验方法。

化学势部分建议明确总量与摩尔量的层级：\(U,S,V\) 是总量，\(\mu\) 是摩尔
Gibbs
自由能。第三问先只写反应计量式，再按``反应物化学势之和等于生成物之和''填入各项。最后分别讨论
\(T\to0\) 与 \(T\to\infty\)，让学生用粒子数比例解释 \(x_p\to0\) 与
\(1/2\)，而不是只做指数极限。

\subsection{九、变式训练}\label{ux4e5dux53d8ux5f0fux8badux7ec3-5}

\begin{enumerate}
\def\labelenumi{\arabic{enumi}.}
\tightlist
\item
  改用 \((V,T)\) 为状态变量，写出 \(S(V,T)\)。
\item
  若质子与中子的参考化学势相差 \(\Delta\mu_0\)，重求 \(x_p\)。
\item
  若电子也按经典理想气体处理，建立三个分压之间的平衡方程。
\end{enumerate}

\textbf{变式解析要点：} 第 1 项用 \(p=nRT/V\) 代回式（5）可得

\[
S(V,T)=S_0+nC_V\ln\frac{T}{T_0}
+nR\ln\frac{V}{V_0},
\]

其中 \(V_0=nRT_0/p_0\)。第 2 项若定义
\(\Delta\mu_0=\mu_{n,0}-\mu_{p,0}\)，则指数中的能量差改为
\(N_A m_ec^2-\Delta\mu_0\)，正负号可由原平衡式重新移项确定。第 3
项需要同时加入电荷守恒与重子数守恒，不能只有一个化学平衡方程；若忽略这些约束，三个分压不会唯一确定。

\newpage

\section{第 7 题}\label{ux7b2c-7-ux9898}

边长为 \(l\) 的最大正方形的四个角上，各连接一个边长为原来 \(1/3\)
的相似正方形；对每个小正方形重复此过程无限次，得到平面分形电阻网络。所有导线单位长度电阻为
\(r\)。垂直纸面向内存在均匀随时间变化的磁场 \(B=kt\)。

（1）计算最大正方形对角线两端之间的等效电阻；

（2）计算最大正方形相邻两顶点之间的等效电阻；

（3）求最外框中间一段长度为 \(l/3\) 的导线中的电流 \(I_c\)；

（4）假设已知总功率可以拆分为四个角上每个相似子网络各自受磁场影响所产生的功率
\(P_1\)，与仅由电流 \(I_c\) 在网络中流动所产生的功率 \(P_2\)
之和，求整个电阻网络单位时间的焦耳热；

（5）证明第（4）问中的总功率确实可以作上述分解。

\begin{figure}[H]
\centering
\includegraphics[width=0.411\linewidth,height=\textheight,keepaspectratio,alt={图7-1　四角自相似正方形电阻网络}]{../figures/question-07.pdf}
\caption*{图7-1　四角自相似正方形电阻网络}
\end{figure}

\subsection{一、题型定位}\label{ux4e00ux9898ux578bux5b9aux4f4d-6}

把每个角上的小网络缩放成四端口，再求外框中段的环流。功率不能随电流直接相加，须对各角部子网络证明交叉项为零。

\subsection{二、知识点}\label{ux4e8cux77e5ux8bc6ux70b9-6}

\begin{enumerate}
\def\labelenumi{\arabic{enumi}.}
\tightlist
\item
  均匀导线电阻与长度成正比。
\item
  自相似网络缩小 \(1/3\) 后，所有等效电阻也缩小 \(1/3\)。
\item
  对称性可把对角端口、相邻端口问题化为有限个未知量。
\item
  法拉第定律 \(\mathcal E=-d\Phi/dt\)。
\item
  线性电阻网络满足叠加原理与互易性。
\item
  总焦耳功率是电流的正定二次型。
\end{enumerate}

\subsection{三、物理图像}\label{ux4e09ux7269ux7406ux56feux50cf-6}

任取一个角上的小网络，把它放大三倍后与整体完全相同。因此``无限深''的信息可以被压缩为有限个端口参数。磁场部分同样自相似：子网络线度缩为
\(1/3\)，面积缩为 \(1/9\)，感应电动势缩为 \(1/9\)，电阻缩为
\(1/3\)，电流与功率分别按 \(1/3\) 和 \(1/27\) 缩放。

\subsection{四、解题入口}\label{ux56dbux89e3ux9898ux5165ux53e3-6}

先把长度 \(l\) 对应的电阻尺度记为
\(R_0=rl\)，以免把``单位长度电阻''与``某段导线的电阻''混淆。再令无量纲端口参数
\(x,y\) 使对角等效电阻为 \(2yR_0\)、相邻等效电阻为
\((x+y)R_0\)，利用自相似性求出 \(x,y\)。磁场部分则把 \(\mathcal E=kl^2\)
作为最大回路的电动势尺度，分别列出顶层环流的等效电动势、等效电阻与功率递推。

\subsection{五、解析}\label{ux4e94ux89e3ux6790-6}

\begin{figure}[H]
\centering
\includegraphics[width=0.891\linewidth,height=\textheight,keepaspectratio,alt={图7-2　端口等效、自相似缩放与中段电流方向}]{../figures/solution-07.pdf}
\caption*{图7-2　端口等效、自相似缩放与中段电流方向}
\end{figure}

图7-2左侧是无源四端口等效，右侧才包含变化磁通的顶层环流。计算端口电阻后，还须保留每个角部的含源电压修正，才能列总回路方程。

\needspace{3\baselineskip}
\textbf{（一）端口等效电阻的自洽方程}

图7-2中，四端口以每个端点串联 \(pR_0\)、内部四边环各边 \(qR_0\)
表示。对应
\(R_{\rm diag}/R_0=2p+q=2y\)、\(R_{\rm adj}/R_0=2p+3q/4=x+y\)，所以
\(p=2x-y,q=4(y-x)\)。这个二参数等效只用于无感应电动势的电阻端口问题。对角接入时，对称支路的等势点合并，得到

\[
2y=\frac23x+\left(\frac13+\frac23y\right). \tag{1}
\]

将括号展开：

\[
2y=\frac23x+\frac13+\frac23y.
\]

两边同减 \(2y/3\)：

\[
\frac43y=\frac23x+\frac13.
\]

乘以 \(3\)：

\[
4y=2x+1,
\]

即

\[
2x-4y=-1. \tag{2}
\]

相邻两点接入时，按电流分流比例合并外框与两个缩小网络，得到

\[
x+y=\frac23x+\frac34\left(\frac13+\frac23y\right). \tag{3}
\]

展开右边：

\[
x+y=\frac23x+\frac14+\frac12y.
\]

移项：

\[
\frac13x+\frac12y=\frac14.
\]

乘以 \(2\)：

\[
\frac23x+y=\frac12. \tag{4}
\]

联立式（2）、（4）。由式（2）有 \(x=2y-1/2\)，代入式（4）：

\[
\frac23\left(2y-\frac12\right)+y=\frac12.
\]

于是

\[
\frac73y=\frac56,
\]

故

\[
y=\frac5{14},\qquad x=\frac3{14}. \tag{5}
\]

因此

\[
\boxed{R_{\rm diag}=2yrl=\frac57rl}, \tag{6}
\]

\[
\boxed{R_{\rm adj}=(x+y)rl=\frac47rl}. \tag{7}
\]

\needspace{3\baselineskip}
\textbf{（二）最外框中段电流}

从本小节开始，记

\[
R_0=rl,
\qquad
\mathcal E=\left|\frac{d\Phi}{dt}\right|=kl^2. \tag{8}
\]

角部子网络的边长为 \(l/3\)，它的对角端口电阻由式（6）按长度缩放为

\[
R_{\rm diag}^{(1)}
=\frac57r\frac l3
=\frac5{21}R_0. \tag{9}
\]

子正方形的面积是最大正方形的 \(1/9\)，因而它自身的感应电动势为

\[
\mathcal E_1=\frac{\mathcal E}{9}. \tag{10}
\]

沿外框绕行时，每个角部子网络都是从一个端点进入、从对角端点流出，因而阻性压降为
\(I_cR_{\rm diag}^{(1)}\)。子网络内的两条对称半路平分式（10）的电动势，故在选定的外框绕行路径上，每个角部对顶层驱动电动势的修正量为

\[
\frac{\mathcal E_1}{2}=\frac{\mathcal E}{18}. \tag{11}
\]

四个角部的修正与最大回路电动势反向，所以真正驱动顶层环流的等效电动势为

\[
\mathcal E_{\rm eff}
=\mathcal E-4\times\frac{\mathcal E}{18}
=\frac79\mathcal E. \tag{12}
\]

顶层环流每经过一边，先经过一段长为 \(l/3\)
的直导线，再经过一个对角端口电阻。因此顶层环流的总电阻为

\[
\begin{aligned}
R_{\rm top}
&=4\left(\frac{R_0}{3}+\frac5{21}R_0\right)\\
&=4R_0\left(\frac7{21}+\frac5{21}\right)\\
&=\frac{16}{7}R_0. 
\end{aligned}
\tag{13}
\]

由欧姆定律，

\[
\begin{aligned}
I_c
&=\frac{\mathcal E_{\rm eff}}{R_{\rm top}}\\
&=\frac{(7/9)\mathcal E}{(16/7)R_0}\\
&=\frac{49}{144}\frac{\mathcal E}{R_0}\\
&=\boxed{\frac{49}{144}\frac{kl}{r}}. 
\end{aligned}
\tag{14}
\]

为了与图 7-2
中的对称半路直接对应，也可记沿选定外框路径每个角部半段的含源电压降为
\(U_1\)，则

\[
U_1=\frac12I_cR_{\rm diag}^{(1)}+\frac14\mathcal E_1
=\frac5{42}R_0I_c+\frac1{36}\mathcal E. \tag{15}
\]

最大回路方程为

\[
8U_1+4I_c\frac{R_0}{3}=\mathcal E. \tag{16}
\]

将式（15）逐项代入：

\[
8\left(\frac5{42}R_0I_c+\frac1{36}\mathcal E\right)
+\frac43R_0I_c=\mathcal E,
\]

\[
\left(\frac{20}{21}+\frac43\right)R_0I_c
=\left(1-\frac29\right)\mathcal E,
\]

\[
\frac{16}{7}R_0I_c=\frac79\mathcal E,
\]

仍得式（14）。

量纲检查：\(kl^2\) 的量纲是伏特，\(rl\) 的量纲是欧姆，故
\(kl^2/(rl)=kl/r\) 的量纲是安培。

\needspace{3\baselineskip}
\textbf{（三）总焦耳功率}

记整个网络的焦耳功率为 \(P(l)\)。对一个边长缩小为 \(l/3\)
的子网络，电动势缩为 \(1/9\)，电阻缩为 \(1/3\)，所以其纯磁场功率相对于
\(P(l)\) 的缩放因子为

\[
\frac{(1/9)^2}{1/3}=\frac1{27}. \tag{17}
\]

四个角部子网络的纯磁场功率之和因而是

\[
P_1=4P\!\left(\frac l3\right)=\frac4{27}P(l). \tag{18}
\]

再计算顶层环流 \(I_c\) 所产生的功率 \(P_2\)。每一边中段的电阻是
\(R_0/3\)，每一角部子网络对环流呈现的对角端口电阻是
\(5R_0/21\)。四组结构串联地通过同一电流 \(I_c\)，故

\[
\begin{aligned}
P_2
&=4I_c^2\left(\frac{R_0}{3}+\frac5{21}R_0\right)\\
&=\frac{16}{7}I_c^2R_0. 
\end{aligned}
\tag{19}
\]

把式（14）先平方：

\[
I_c^2
=\left(\frac{49}{144}\right)^2\frac{\mathcal E^2}{R_0^2}
=\frac{2401}{20736}\frac{\mathcal E^2}{R_0^2}. \tag{20}
\]

代入式（19）：

\[
\begin{aligned}
P_2
&=\frac{16}{7}\cdot\frac{2401}{20736}
\frac{\mathcal E^2}{R_0}\\
&=\frac{16\times343}{20736}\frac{\mathcal E^2}{R_0}\\
&=\frac{343}{1296}\frac{\mathcal E^2}{R_0}. 
\end{aligned}
\tag{21}
\]

由第（4）问给出的功率分解，

\[
P(l)=P_1+P_2
=\frac4{27}P(l)+\frac{343}{1296}\frac{\mathcal E^2}{R_0}. \tag{22}
\]

把递归项移到左边：

\[
\frac{23}{27}P(l)
=\frac{343}{1296}\frac{\mathcal E^2}{R_0}. \tag{23}
\]

因为 \(1296/27=48\)，所以

\[
\begin{aligned}
P(l)
&=\frac{343}{1296}\cdot\frac{27}{23}
\frac{\mathcal E^2}{R_0}\\
&=\frac{343}{48\times23}\frac{\mathcal E^2}{R_0}\\
&=\frac{343}{1104}\frac{\mathcal E^2}{R_0}. 
\end{aligned}
\tag{24}
\]

最后代回 \(\mathcal E=kl^2\) 与 \(R_0=rl\)：

\[
\boxed{P(l)=\frac{343}{1104}\frac{k^2l^3}{r}}. \tag{25}
\]

量纲上，\(\mathcal E^2/R_0\) 的单位是
\(\mathrm{V^2}/\Omega=\mathrm W\)；等价地，\(k^2l^3/r\) 也是功率。

\needspace{3\baselineskip}
\textbf{（四）功率可加性的证明}

先对每一个角部子网络截断到有限层，在该子网络上分解电流。设关闭子网络内感应源、只由其端口电流
\(I_c\) 产生的各支路电流为
\(\bm i_I\)，该子网络外端口断开、只由其内部变化磁场产生的支路电流为
\(\bm i_B\)。由线性网络的叠加原理，总电流为

\[
\bm i=\bm i_I+\bm i_B. \tag{26}
\]

设 \(\mathbf R\) 为以各支路电阻为对角元素的对角矩阵。总焦耳功率展开为

\[
\begin{aligned}
P
&=(\bm i_I+\bm i_B)^{\mathsf T}
\mathbf R(\bm i_I+\bm i_B)\\
&=P_I+P_B
+2\bm i_I^{\mathsf T}\mathbf R\bm i_B. 
\end{aligned}
\tag{27}
\]

因此问题的核心是证明交叉项为零。在纯端口激励问题中，内部没有感应电动势，第
\(j\) 条支路的阻性电压为

\[
v_{I,j}=R_ji_{I,j}. \tag{28}
\]

将式（28）与纯磁场解的支路电流相乘并对这一个角部子网络求和：

\[
\sum_jR_ji_{I,j}i_{B,j}
=\sum_jv_{I,j}i_{B,j}. \tag{29}
\]

右边可按结点重新归并。设纯端口激励下第 \(n\) 个结点电势为
\(\varphi_{I,n}\)，则

\[
\sum_jv_{I,j}i_{B,j}
=\sum_n\varphi_{I,n}
\left(\sum_{j\in n}i_{B,j}\right). \tag{30}
\]

纯磁场解没有从外部端口注入的净电流，而每一个内部结点也满足电流守恒，所以括号中的和对每个结点都是零。因此

\[
\bm i_I^{\mathsf T}\mathbf R\bm i_B=0. \tag{31}
\]

代回式（27），得

\[
\boxed{P=P_I+P_B}. \tag{32}
\]

对任意有限截断层数上式均成立；同时，子网络功率的缩放系数为
\(4/27<1\)，故当截断层数趋于无穷时功率级数收敛，式（32）仍然成立。

\needspace{3\baselineskip}
\textbf{自相似缩放与功率正交性的复核}

端口参数 \(x,y\)
本身必须是无量纲数，因为所有等效电阻都被写成``无量纲系数乘
\(rl\)''。当子网络的边长由 \(l\) 缩为 \(l/3\) 时，单位长度电阻 \(r\)
不变，任一对应导线电阻都乘 \(1/3\)，所以子网络端口电阻也整体乘
\(1/3\)。若自洽方程中出现 \(x/9\) 或
\(3x\)，便说明把长度缩放与面积缩放混淆了。

解出的两个等效电阻满足

\[
R_{\rm diag}=\frac57rl,
\qquad
R_{\rm adj}=\frac47rl.
\]

二者均小于沿外框直接走一条或两条边的电阻，符合大量并联支路会降低等效电阻的直觉；同时对角端口路径整体较长，因此
\(R_{\rm diag}>R_{\rm adj}\)。这两个不等式是无需重算网络即可完成的数量级检查。

磁场部分要分开记录长度、面积、电动势、电阻、电流和功率的缩放。边长乘
\(1/3\)，面积乘 \(1/9\)，故感应电动势

\[
\mathcal E\sim k l^2
\]

乘 \(1/9\)；电阻乘 \(1/3\)，所以对应电流乘

\[
\frac{1/9}{1/3}=\frac13.
\]

功率既可按 \(I^2R\) 计算，也可按 \(\mathcal E^2/R\)
计算，两种方法都给出单个子网络的缩放因子

\[
\left(\frac13\right)^2\left(\frac13\right)=\frac1{27}.
\]

四个子网络总系数是 \(4/27<1\)，因而功率递推级数收敛。若缩放比一般为
\(s\)，单个子网络功率按 \(s^3\) 缩放；有四个子网络时，有限总功率至少要求
\(4s^3<1\)。

最后解释交叉项为何恰好为零。对纯端口激励解，内部每个独立回路的外加电动势为零；对纯磁场激励解，外端口的净注入电流为零。把第一组回路电压方程乘第二组回路电流并求和，再把第二组结点电流方程与第一组结点电势配对，边界项均消失，得到

\[
\sum_j R_j i_{I,j}i_{B,j}=0.
\]

因此

\[
\sum_jR_j(i_{I,j}+i_{B,j})^2
=\sum_jR_ji_{I,j}^2+
\sum_jR_ji_{B,j}^2.
\]

这里可相加的是这两组特殊激励的总功率，并不是说任意两组电流叠加后功率都可直接相加。只有交叉二次型为零时才成立。

\needspace{3\baselineskip}
\textbf{把局部正交分解拼成全网络}

式（26）---（32）的证明应分别应用于四个角部。纯磁场分量在该子网络每个外端口的净电流为零；纯端口分量在其内部没有非保守源，因而
\(R_ji_{I,j}\) 可写为结点势差。四条顶层连接段没有子网络纯磁场分量，只含
\(I_c\)，功率直接为
\(4I_c^2R_0/3\)。把四个局部正交分解和这四段相加，才得到式（22）。不能把整个闭合网络凭空说成一个外端口供电、无内部感应源的网络。

若 \(k>0\) 且磁场向纸内增长，\(I_c\)
绕外框逆时针；图7-2的顶边箭头应向左。相似图只画有限层，点线表示继续递归；相互重合的导线为同一条导线，不能重复计电阻。

\needspace{3\baselineskip}
\textbf{用有限网络核对无限极限}

取
\(r=l=k=1\)，从单个外框开始，按原图四角不断缩小三分之一。把所有重合线段合并后，在实际结点列基尔霍夫方程；一条连续导线不能因同时属于两层正方形而被重复并联。

有限层得到下列值，层数0表示只有外框：

{\def\LTcaptype{none} % do not increment counter
\begin{longtable}[]{@{}lllll@{}}
\toprule\noalign{}
层数 & 对角电阻 & 相邻电阻 & 外框中段电流 & 总功率 \\
\midrule\noalign{}
\endhead
\bottomrule\noalign{}
\endlastfoot
0 & 1.000000 & 0.750000 & 0.250000 & 0.250000 \\
1 & 0.833333 & 0.666667 & 0.291667 & 0.263889 \\
2 & 0.777778 & 0.625000 & 0.318182 & 0.286569 \\
3 & 0.750000 & 0.601852 & 0.328125 & 0.297663 \\
5 & 0.725823 & 0.581276 & 0.336299 & 0.306508 \\
无限层 & \(5/7\) & \(4/7\) & \(49/144\) & \(343/1104\) \\
\end{longtable}
}

有限值逐层向正文的四个极限收敛。这个检查直接使用几何导线图、每段长度和感应电动势，没有把无限层答案先塞进每个角部。

每条有向线段从 \((x_a,y_a)\) 到 \((x_b,y_b)\) 的感应电动势可取
\((k/2)(x_ay_b-y_ax_b)\)，其各回路和等于面积乘k。改变感应电场的梯度规范只改变结点势，不改变电流、热功率与上述极限。

\subsection{六、关键洞察}\label{ux516dux5173ux952eux6d1eux5bdf-6}

解析图左侧是无源四端口等效，右侧才包含变化磁通的顶层环流。计算端口电阻后，还须保留每个角部的含源电压修正，才能列总回路方程。

第五问可以看成线性网络中的``正交分解''。端口电流激励与闭合回路磁通激励占据互补的边界条件：前者由外部注入电流，后者在端口无净注入。它们在电阻加权内积

\[
\langle\bm i_1,\bm i_2\rangle_R
=\bm i_1^{\mathsf T}\mathbf R\bm i_2
\]

下正交，所以能量二次型没有交叉项。这一结论不是普通叠加原理的附带结果，而是网络互易性与边界条件共同给出的更强性质。

\subsection{七、易错点}\label{ux4e03ux6613ux9519ux70b9-6}

\begin{enumerate}
\def\labelenumi{\arabic{enumi}.}
\tightlist
\item
  缩小 \(1/3\) 后忘记等效电阻也缩小 \(1/3\)。
\item
  把相邻等效电阻误写为 \(xrl\)，漏掉另一部分 \(yrl\)。
\item
  把单位长度电阻 \(r\) 直接当成某段导线的电阻。正确的电阻尺度是
  \(R_0=rl\)，长 \(l/3\) 的导线电阻是 \(R_0/3\)。
\item
  感应电动势方向没有统一，把角部修正错写成与最大回路电动势同向。
\item
  功率递推把四个子网络写成 \(4/9\) 或 \(4/81\)，正确为 \(4/27\)。
\item
  直接声称功率可叠加；电流可叠加不自动意味着功率可叠加，必须证明交叉项为零。
\end{enumerate}

批改时可把答案拆成三层：端口自洽方程、磁场回路电流、功率递推。每层都要出现正确的缩放因子。因为
\(r\) 是单位长度电阻，总功率最后应正比于 \(k^2l^3/r\)；若得到
\(k^2l^4/r\)，就是把 \(r\)
误当成电阻而漏掉了长度因子。第五问必须至少写出二次型交叉项，并说明它因端口与回路边界条件正交而消失。

\subsection{八、课堂讲解建议}\label{ux516bux8bfeux5802ux8bb2ux89e3ux5efaux8bae-6}

先用有限两层网络让学生观察端口等效，再让层数趋于无穷。功率部分重点让学生自己推导缩放因子
\((1/9)^2/(1/3)\)。最后用二次型展开说明为何第五问不是普通叠加原理的直接结论。

建议先把无限图形截断到一层、两层，用同一对端口分别测量相邻与对角等效电阻，使学生意识到至少需要两个端口参数。建立自洽方程后，每个
\(1/3\)
都要求说清来源：是线长、端口电阻，还是电流分配。这样可以避免面积缩放因子过早混入电阻方程。

磁场部分可画一个``缩放账本''：长度 \(1/3\)、面积 \(1/9\)、电动势
\(1/9\)、电阻 \(1/3\)、电流 \(1/3\)、功率
\(1/27\)。最后让学生把总功率展开成三项，并回答``电流能叠加，为什么功率通常不能叠加''。只有他说出交叉项以及本题交叉项为零的网络条件，才算真正理解第五问。

\subsection{九、变式训练}\label{ux4e5dux53d8ux5f0fux8badux7ec3-6}

\begin{enumerate}
\def\labelenumi{\arabic{enumi}.}
\tightlist
\item
  把缩放比 \(1/3\) 改为一般 \(s\)，建立端口自洽方程。
\item
  若磁场满足 \(B=B_0\sin\omega t\)，求功率的时间平均值。
\item
  在对角端口接入理想电压源，求总功率中交叉项成立的条件。
\end{enumerate}

\textbf{变式解析要点：} 一般缩放比 \(s\) 下，子网络电阻按 \(s\)
缩放，磁场电动势按 \(s^2\) 缩放，故电流按 \(s\)、功率按 \(s^3\)
缩放。若仍在四角复制，功率递推的收敛条件是 \(4s^3<1\)。当
\(B=B_0\sin\omega t\) 时，\(k\) 应换成瞬时
\(B_0\omega\cos\omega t\)；正文中功率与 \(k^2\) 成正比，所以周期平均再乘
\(1/2\)。端口改接电压源时，应重新检查纯磁场解的端口电流是否为零；只有外端口边界项仍消失，交叉功率才保持为零。

\newpage

\section{第 8 题}\label{ux7b2c-8-ux9898}

交流电源 \(U(t)=U_0\cos\omega t\)
接在``日''字形理想变压器的左侧线圈上。三根铁芯柱上的线圈匝数为
\(N_1,N_2,N_3\)；输入回路串联电阻为 \(r_1'\)，右侧负载含
\(r_1,r_2,r_3\)。忽略漏磁、铁损和线圈电阻。求稳定工作时 \(r_2\)
上的平均功率 \(P_2\)，并在 \(N_1:N_2:N_3=4:2:1\) 时化简。

\begin{figure}[H]
\centering
\includegraphics[width=0.84\linewidth,height=\textheight,keepaspectratio,alt={图8-1　日字形铁芯、三线圈与负载连接}]{../figures/question-08.pdf}
\caption*{图8-1　日字形铁芯、三线圈与负载连接}
\end{figure}

\subsection{一、题型定位}\label{ux4e00ux9898ux578bux5b9aux4f4d-7}

三柱磁通先分流，再通过安匝平衡约束电流。公共电阻流过两个次级电流之和，平方折算后出现交叉项；最后用峰值转有效值求平均功率。

\subsection{二、知识点}\label{ux4e8cux77e5ux8bc6ux70b9-7}

\begin{enumerate}
\def\labelenumi{\arabic{enumi}.}
\tightlist
\item
  铁芯磁导率趋于无穷时，铁芯内磁场强度 \(H\to0\)。
\item
  安培环路定律给出磁动势平衡。
\item
  法拉第定律给出绕组电压 \(u_j=N_j\dot\Phi_j\)。
\item
  磁通在铁芯分叉处满足连续性 \(\Phi_1=\Phi_2+\Phi_3\)。
\item
  正弦电流若峰值为 \(I_0\)，则电阻平均功率为 \(I_0^2R/2\)。
\end{enumerate}

\subsection{三、物理图像}\label{ux4e09ux7269ux7406ux56feux50cf-7}

左柱磁通进入上横轭后分成中柱磁通和右柱磁通。理想铁芯不需要有限磁场强度维持磁通，因此各窗口上的绕组磁动势必须互相抵消。右侧中点电阻
\(r_1\) 同时承受两个次级支路电流之和，因而在折算到原边时出现交叉耦合项。

\subsection{四、解题入口}\label{ux56dbux89e3ux9898ux5165ux53e3-7}

先用安培环路定律把次级电流写成原边电流的匝数比倍数；再写三个绕组的电压方程。利用
\(\dot\Phi_1=\dot\Phi_2+\dot\Phi_3\) 消去磁通，只留下原边电流，最后求
\(r_2\) 的电流有效值。

\subsection{五、解析}\label{ux4e94ux89e3ux6790-7}

\begin{figure}[H]
\centering
\includegraphics[width=0.727\linewidth,height=\textheight,keepaspectratio,alt={图8-2　三柱磁通方向、绕组电流与负载支路}]{../figures/solution-08.pdf}
\caption*{图8-2　三柱磁通方向、绕组电流与负载支路}
\end{figure}

图8-2按三个磁柱的通量和两个次级共结点列约束，公共电阻流过两次级电流之和。阻值折算的交叉项来自这个实际连接。

\needspace{3\baselineskip}
\textbf{（一）磁通与磁动势关系}

按图 8-2 取磁通正方向。分叉处磁通连续：

\[
\Phi_1=\Phi_2+\Phi_3. \tag{1}
\]

理想铁芯中 \(H\to0\)。分别沿两个窗口应用安培环路定律，按选定方向可写成

\[
N_1I_1=N_2I_2=N_3I_3. \tag{2}
\]

定义匝数折算系数

\[
a=\frac{N_1}{N_2},\qquad b=\frac{N_1}{N_3}. \tag{3}
\]

于是

\[
I_2=aI_1,\qquad I_3=bI_1. \tag{4}
\]

公共电阻 \(r_1\) 中电流为

\[
I_a=I_2+I_3=(a+b)I_1. \tag{5}
\]

\needspace{3\baselineskip}
\textbf{（二）三个绕组的电压方程}

以下均用同一相位的峰值相量表示。原边满足

\[
N_1\dot\Phi_1=U_0-I_1r_1'. \tag{6}
\]

中柱次级回路：

\[
N_2\dot\Phi_2=r_2I_2+r_1(I_2+I_3). \tag{7}
\]

右柱次级回路：

\[
N_3\dot\Phi_3=r_3I_3+r_1(I_2+I_3). \tag{8}
\]

由式（1）：

\[
N_1\dot\Phi_1
=\frac{N_1}{N_2}N_2\dot\Phi_2
+\frac{N_1}{N_3}N_3\dot\Phi_3. \tag{9}
\]

将式（7）、（8）代入式（9），再用式（4）、（5）：

\[
\begin{aligned}
N_1\dot\Phi_1
={}&a\left[r_2(aI_1)+r_1(a+b)I_1\right]\\
&+b\left[r_3(bI_1)+r_1(a+b)I_1\right].
\end{aligned}
\]

提取 \(I_1\)：

\[
N_1\dot\Phi_1
=I_1\left[a^2r_2+b^2r_3+(a+b)^2r_1\right]. \tag{10}
\]

联立式（6）、（10）：

\[
U_0-I_1r_1'=I_1\left[a^2r_2+b^2r_3+(a+b)^2r_1\right].
\]

因此原边电流峰值为

\[
I_1=\frac{U_0}{r_1'+a^2r_2+b^2r_3+(a+b)^2r_1}. \tag{11}
\]

记

\[
D=r_1'+a^2r_2+b^2r_3+(a+b)^2r_1. \tag{12}
\]

则 \(r_2\) 中的电流峰值为

\[
I_2=aI_1=\frac{aU_0}{D}. \tag{13}
\]

\needspace{3\baselineskip}
\textbf{（三）\(r_2\) 的平均功率}

正弦电流有效值为峰值的 \(1/\sqrt2\)，故

\[
P_2=\left(\frac{I_2}{\sqrt2}\right)^2r_2.
\]

代入式（13）：

\[
\boxed{P_2=\frac12r_2a^2\left(\frac{U_0}{D}\right)^2}. \tag{14}
\]

写回匝数：

\[
\boxed{
P_2=\frac12r_2\left(\frac{N_1}{N_2}\right)^2
\left[
\frac{U_0}{
r_1'+\left(\frac{N_1}{N_2}\right)^2r_2
+\left(\frac{N_1}{N_3}\right)^2r_3
+\left(\frac{N_1}{N_2}+\frac{N_1}{N_3}\right)^2r_1}
\right]^2}. \tag{15}
\]

\needspace{3\baselineskip}
\textbf{（四）匝数比 \(4:2:1\)}

此时

\[
a=2,\qquad b=4,\qquad a+b=6.
\]

所以

\[
D=r_1'+4r_2+16r_3+36r_1. \tag{16}
\]

且式（14）中的系数为 \((1/2)a^2=2\)，最终

\[
\boxed{P_2=\frac{2r_2U_0^2}{\left(r_1'+36r_1+4r_2+16r_3\right)^2}}. \tag{17}
\]

若所有负载趋于无穷大，则 \(I_1,I_2\to0\)，功率趋于零；若
\(r_2\to0\)，式（17）也给出 \(P_2\to0\)，两个极限均合理。

\needspace{3\baselineskip}
\textbf{磁通拓扑与折算阻抗的复核}

``日''字形铁芯的三个柱并不是三个互不相关的理想变压器。先选定各柱磁通正方向，由结点处磁通连续性得到

\[
\Phi_1=\Phi_2+\Phi_3.
\]

对时间求导并乘相应匝数，才得到各线圈感应电动势之间的关系。若某个次级线圈反接，改变的是电动势与电流的参考方向，公共支路中的和会相应变成差；不能只在最后公式随意改号。

折算到原边时，每个次级电流先按安匝平衡换算：

\[
I_2=\frac{N_1}{N_2}I_1=aI_1,
\qquad
I_3=\frac{N_1}{N_3}I_1=bI_1.
\]

因此单独流过 \(r_2\)、\(r_3\) 的功率分别为
\(a^2I_1^2r_2\)、\(b^2I_1^2r_3\)。公共电阻 \(r_1\) 中流过的是
\(I_2+I_3=(a+b)I_1\)，故其折算电阻必为 \((a+b)^2r_1\)，不是
\((a^2+b^2)r_1\)。从功率等价

\[
I_1^2R_{\rm eq}=I_2^2r_2+I_3^2r_3+(I_2+I_3)^2r_1
\]

也能不依赖电路图直接复原分母 \(D\)。

题给 \(U_0\) 是电压峰值，故原边电流峰值为 \(I_{1m}=U_0/D\)，\(r_2\)
电流峰值为 \(aU_0/D\)。正弦稳态平均功率应使用有效值：

\[
P_2=r_2\left(\frac{aU_0}{\sqrt2D}\right)^2,
\]

这正好给出式（14）中的 \(1/2\)。量纲上，\(r_2U_0^2/D^2\) 为
\(\Omega\cdot\mathrm{V^2}/\Omega^2=\mathrm W\)。

最后做两个互补极限：\(r_2\to0\)
时该支路虽有电流但电阻为零，\(P_2\to0\)；\(r_2\to\infty\)
时电流被抑制，亦有 \(P_2\to0\)。所以在其余参数固定时必存在有限的最佳
\(r_2\)。这也说明若推导所得功率随 \(r_2\)
单调增大，分母中的折算项一定有遗漏。

\needspace{3\baselineskip}
\textbf{原图符号与最大功率的补充}

原扫描中输入串联电阻和右侧公共电阻均标为 \(r_1\)。整理时以 \(r_1'\)
区分输入电阻；若原题要求两处取同一个值，则式（17）应写成

\[
P_2=\frac{2r_2U_0^2}{(37r_1+4r_2+16r_3)^2}.
\]

图8-2分别标出两电阻的位置及次级共同结点，以免仅凭撇号猜连接。设
\(C=r_1'+b^2r_3+(a+b)^2r_1>0\)，最大功率条件为 \(r_2=C/a^2\)，相应
\(P_{2,\max}=U_0^2/(8C)\)。源输入平均功率 \(U_0I_1/2\)
等于全部电阻平均耗散，另给一个总体检查。

\subsection{六、关键洞察}\label{ux516dux5173ux952eux6d1eux5bdf-7}

解析图按三个磁柱的通量和两个次级共结点列约束，公共电阻流过两次级电流之和。阻值折算的交叉项来自这个实际连接。

\subsection{七、易错点}\label{ux4e03ux6613ux9519ux70b9-7}

\begin{enumerate}
\def\labelenumi{\arabic{enumi}.}
\tightlist
\item
  把 \(\Phi_1=\Phi_2+\Phi_3\) 写成三个磁通相等。
\item
  漏掉公共电阻电流 \(I_2+I_3\)。
\item
  折算阻抗时只乘一次匝数比；电阻折算必须乘匝数比平方。
\item
  用峰值电流直接计算平均功率，漏掉因子 \(1/2\)。
\item
  匝数比代入后把 \((2+4)^2\) 错算成 \(2^2+4^2\)。
\end{enumerate}

验算应按``磁通连续---安匝平衡---公共支路电流---阻抗平方折算---峰值转有效值''的顺序逐项打勾。最终分母中的每一项都必须具有电阻量纲且为非负；把任一负载开路时，对应支路功率应趋于零。若线圈反接，只有带方向的公共项改变，单支路折算电阻仍为正。

\subsection{八、课堂讲解建议}\label{ux516bux8bfeux5802ux8bb2ux89e3ux5efaux8bae-7}

先让学生在铁芯上画箭头并写式（1），再沿左右两个窗口写磁动势平衡。负载部分用结点电流图解释公共电阻为何流过
\(I_2+I_3\)，最后再讨论阻抗折算。

教学时可先把所有匝数比都设为
\(1\)。此时公共支路电流明显是两支路电流之和，等效分母中应出现
\(4r_1\)；再恢复一般匝数比，学生更容易接受
\((a+b)^2r_1\)。若一上来就写折算公式，常会把公共电阻误拆成两个独立负载。

计算功率前要在黑板上注明``题给 \(U_0\)
为峰值''。让学生分别用峰值和有效值写一遍电流，观察 \(1/2\)
从何而来。结尾可固定其余电阻，讨论 \(r_2\to0\)、\(r_2\to\infty\)
两端功率都为零，由连续性判断中间必有最大值，为变式中的最大功率传输问题作铺垫。

\subsection{九、变式训练}\label{ux4e5dux53d8ux5f0fux8badux7ec3-7}

\begin{enumerate}
\def\labelenumi{\arabic{enumi}.}
\tightlist
\item
  若右柱线圈反接，重写公共支路电流与 \(P_2\)。
\item
  若中柱与右柱存在有限磁阻，建立磁路方程。
\item
  保持 \(r_1',r_1,r_3\) 不变，求使 \(P_2\) 最大的 \(r_2\)。
\end{enumerate}

\textbf{变式解析要点：} 记

\[
C=r_1'+b^2r_3+(a+b)^2r_1,
\]

则

\[
P_2=\frac12a^2U_0^2\frac{r_2}{(C+a^2r_2)^2}.
\]

对 \(r_2\) 求导：

\[
\frac{d}{dr_2}\frac{r_2}{(C+a^2r_2)^2}
=\frac{C-a^2r_2}{(C+a^2r_2)^3}.
\]

因此最大值条件为

\[
a^2r_2=C,
\]

即负载折算到原边后的电阻等于其余串联等效电阻。这正是最大功率传输定理在本题多柱铁芯拓扑中的表现。若右柱线圈反接，只需先把公共支路的
\((a+b)\) 改为 \((a-b)\)，再重复同一推导。
\end{document}
